\documentclass{article} % For LaTeX2e
\usepackage{iclr2026_conference,times}

\usepackage[most]{tcolorbox}
\usepackage{xcolor}
\definecolor{abstractbg}{RGB}{235, 245, 255} 
\definecolor{metafg}{HTML}{1C2B33}

\tcbset{
  enhanced,
  colback=abstractbg,      % ✅ 背景色
  colframe=metafg,     % ✅ 边框颜色同背景
  boxrule=0pt,             % ✅ 去掉边线
  frame hidden,            % ✅ 彻底隐藏边框（双保险）
  arc=8pt,                 % ✅ 圆角更大（可调 10pt、12pt）
  left=0pt,right=0pt,top=6pt,bottom=6pt,
  before skip=8pt, after skip=12pt,
  % 可选：长摘要自动分页
  % breakable,
}

\newenvironment{boxedabstract}
{\begin{tcolorbox}\begin{abstract}}
{\end{abstract}\end{tcolorbox}}


% Optional math commands from https://github.com/goodfeli/dlbook_notation.
\input{math_commands.tex}


\usepackage{multirow}
\usepackage[utf8]{inputenc} % allow utf-8 input
\usepackage{url}            % simple URL typesetting
\usepackage{tabularx} 
\usepackage{array}
\renewcommand{\arraystretch}{1.1}
\usepackage[T1]{fontenc}    % use 8-bit T1 fonts
\usepackage{booktabs}       % professional-quality tables
\usepackage{amsfonts}       % blackboard math symbols
\usepackage{nicefrac}       % compact symbols for 1/2, etc.
\usepackage{microtype}      % microtypography
\usepackage{xcolor}         % colors

\usepackage{multirow}
\usepackage[table,xcdraw]{xcolor}

\definecolor{metablue}{HTML}{0064E0}

\usepackage{caption}
\usepackage{subcaption}
\usepackage{natbib}
\usepackage{graphicx}
\usepackage{enumitem}
\usepackage{makecell}
\usepackage{amsmath}
\usepackage{amsthm}
\usepackage{thmtools}
\usepackage{thm-restate}
\usepackage{algorithm}
\usepackage{algorithmic}
\usepackage{bbm}
\usepackage{bm}
\usepackage{bbding} 

% new packages
\usepackage{graphicx}
\usepackage{wrapfig}
\usepackage{multirow}
\usepackage{tcolorbox}
\usepackage{adjustbox} 
\usepackage{xcolor}
\usepackage{float}
\usepackage[normalem]{ulem}
\useunder{\uline}{\ul}{}

\RequirePackage{hyperref}
\RequirePackage[noabbrev,nameinlink]{cleveref}
\hypersetup{
  colorlinks,
  linkcolor=metablue,
  citecolor=metablue,
  urlcolor=metablue
}

\newcommand\norm[1]{\left\lVert#1\right\rVert}
\newcommand{\Lapl}{\mathbf{\mathop{\mathcal{L}}}}
\newcommand{\lapl}{\mathcal{L}}
\newcommand{\Trans}[1]{{#1}^{\top}}
\newcommand{\Trace}[1]{tr\left({#1}\right)}
\newcommand{\Bracs}[1]{\left({#1}\right)}
\newcommand{\Mat}[1]{\mathbf{#1}}
\newcommand{\MatS}[3]{\mathbf{#1}^{#2}_{#3}}
\newcommand{\Space}[1]{\mathbb{#1}}
\newcommand{\Set}[1]{\mathcal{#1}}
\newcommand{\vectornorm}[1]{\left|\left|#1\right|\right|}
\newcommand{\bpi}{\boldsymbol{\pi}}
\newcommand{\BlockMat}[2]{\left[\begin{matrix}#1\\#2\end{matrix}\right]}
\newcommand{\BlockMatSquare}[4]{\left[\begin{matrix}#1 & #2\\#3 & #4\end{matrix}\right]}
\newcommand{\ParaTitle}[1]{\noindent\textbf{#1}}
\newcommand\myeq{\mathrel{\stackrel{\makebox[0pt]{\mbox{\normalfont\tiny def}}}{=}}}
\newcommand\eqc{\mathrel{\overset{\makebox[0pt]{\mbox{\normalfont\tiny\sffamily c}}}{=}}}


\newcommand{\ie}{\emph{i.e., }}
\newcommand{\eg}{\emph{e.g., }}
\newcommand{\etal}{\emph{et al.}}
\newcommand{\st}{\emph{s.t. }}
\newcommand{\etc}{\emph{etc.}}
\newcommand{\wrt}{\emph{w.r.t.}~}
\newcommand{\cf}{\emph{cf. }}
\newcommand{\aka}{\emph{aka. }}
\newcommand{\token}[1]{\texttt{\textless#1\textgreater}}

\newcommand{\email}[1]{\href{mailto:#1}{\texttt{#1}}}

\newcommand{\wx}[1]{{\color{blue}{#1}}}
\newcommand{\za}[1]{{\color{black}{#1}}}
\newcommand{\kxy}[1]{{\color{black}{#1}}}
\newcommand{\tjf}[1]{{\color{black}{#1}}}
\newcommand{\wjc}[1]{{\color{black}{#1}}}
\newcommand{\slh}[1]{{\color{black}{#1}}}
\usepackage{fontawesome} 

%% 页眉
%% 页眉

\usepackage{graphicx}
\usepackage{fancyhdr}

\newcommand{\HeaderIconHeight}{12.3pt}
\newcommand{\CornerIconHeight}{27pt} % 角标尺寸（可改 14–18pt）
% 在导言区新增
\newcommand{\NextIconHeight}{13pt}
\newcommand{\AlphaIconHeight}{12pt}
\newcommand{\LdsIconHeight}{14.5pt}

\newcommand{\CornerRaise}{1.2pt}       % 角标上抬
\newcommand{\CornerGapA}{4pt}        % next 与 alphalab 间距
\newcommand{\CornerGapB}{4pt}        % alphalab 与 lds 间距

\setlength{\headheight}{\HeaderIconHeight}
\pagestyle{fancy}
\addtolength{\headsep}{-27pt}
\fancyhf{}

% 图标放页眉（偶数页左、奇数页右）
\fancyhead[LO,RE]{%
  \raisebox{0pt}[0pt][0pt]{
\includegraphics[height=\HeaderIconHeight]{figures/logov6.png}}%
}

% % 外侧角：next.jpg（最外侧）— alphalab.png — lds.png（最内侧）
% \fancyhead[RO,LE]{%
%   \raisebox{\CornerRaise}[0pt][0pt]{%
%     
\includegraphics[height=\LdsIconHeight]{figures/lds.png}%
%     \hspace{\CornerGapA}%
%     
\includegraphics[height=\AlphaIconHeight]{figures/alphalab.jpg}%
%     \hspace{\CornerGapB}%
%     
\includegraphics[height=\NextIconHeight]{figures/next.jpg}%
%   }%
% }

% 打开页眉横线（调整粗细）
\renewcommand{\headrulewidth}{1pt}
\newcommand{\HeadRuleRaise}{0pt} % 想更高就调大些
\makeatletter
\renewcommand{\headrule}{%
  \vspace*{-\HeadRuleRaise}% 往上抬
  \hrule height \headrulewidth \@width\headwidth
  \vspace*{-\headrulewidth}% 保持竖向占位
}
\makeatother

\fancyfoot[C]{\thepage}

% 首页（plain）同款设置 + 横线
\fancypagestyle{plain}{%
  \fancyhf{}
  \fancyhead[LO,RE]{%
    \raisebox{0pt}[0pt][0pt]{
\includegraphics[height=\HeaderIconHeight]{figures/logov6.png}}%
  }
  % % plain 里也加角落：next — alphalab — lds
  % \fancyhead[RO,LE]{%
  %   \raisebox{\CornerRaise}[0pt][0pt]{%
  %     
\includegraphics[height=\LdsIconHeight]{figures/lds.png}%
  %     \hspace{\CornerGapA}%
  %     
\includegraphics[height=\AlphaIconHeight]{figures/alphalab.png}%
  %     \hspace{\CornerGapB}%
  %     
\includegraphics[height=\NextIconHeight]{figures/next.jpg}%
  %   }%
  % }

  \renewcommand{\headrulewidth}{0.4pt}%
  \makeatletter
  \renewcommand{\headrule}{%
    \vspace*{-\HeadRuleRaise}%
    \hrule height \headrulewidth \@width\headwidth
    \vspace*{-\headrulewidth}%
  }
  \makeatother

  \fancyfoot[C]{\thepage}
}


\title{MiniOneRec: An Open-Source Framework for \\Scaling Generative Recommendation}

% Authors must not appear in the submitted version. They should be hidden
% as long as the \iclrfinalcopy macro remains commented out below.
% Non-anonymous submissions will be rejected without review.

\author{\textbf{Xiaoyu Kong}$^{\star}$,\quad\textbf{Leheng Sheng}$^{\circ}$,\quad\textbf{Junfei Tan}$^{\star}$,\quad\textbf{Yuxin Chen}$^{\circ}$,\\
\textbf{Jiancan Wu}$^{\star}$,\quad\textbf{An Zhang}$^{\star}$,\quad\textbf{Xiang Wang}$^{\star}$,\quad\textbf{Xiangnan He}$^{\star}$\\
  $^\star$University of Science and Technology of China~~
  $^\circ$National University of Singapore\\
  \email{\{kongxy,sober\_clever\}@mail.ustc.edu.com},\\
  %\email{\{wujcan,an.zhang3.14\}@gmail.com},\\
  \email{\{leheng.sheng,yuxin.chen\}@u.nus.edu},\\
  \email{\{wujcan,an.zhang3.14,xiangwang1223,xiangnanhe\}@gmail.com}
  % ,\quad
  % \email{xiangnanhe@gmail.com}
}


% The \author macro works with any number of authors. There are two commands
% used to separate the names and addresses of multiple authors: \And and \AND.
%
% Using \And between authors leaves it to \LaTeX{} to determine where to break
% the lines. Using \AND forces a linebreak at that point. So, if \LaTeX{}
% puts 3 of 4 authors names on the first line, and the last on the second
% line, try using \AND instead of \And before the third author name.

\newcommand{\fix}{\marginpar{FIX}}
\newcommand{\new}{\marginpar{NEW}}

\iclrfinalcopy % Uncomment for camera-ready version, but NOT for submission.
\begin{document}


\maketitle


\begin{boxedabstract}

The recent success of large language models (LLMs) has renewed interest in whether recommender systems can achieve similar scaling benefits. Conventional recommenders, dominated by massive embedding tables, tend to plateau as embedding dimensions grow. In contrast, the emerging generative paradigm replaces embeddings with compact Semantic ID (SID) sequences produced by autoregressive Transformers. Yet most industrial deployments remain proprietary, leaving two fundamental questions open: (1) Do the expected scaling laws hold on public benchmarks? (2) What is the minimal post-training recipe that enables competitive performance?

We present \textbf{MiniOneRec}, to the best of our knowledge, the first fully open-source generative recommendation framework, which provides an end-to-end workflow spanning SID construction, supervised fine-tuning, and recommendation-oriented reinforcement learning. We generate SIDs via a Residual Quantized VAE and post-train Qwen backbones ranging from 0.5B to 7B parameters on the Amazon Review dataset. Our experiments reveal a consistent downward trend in both training and evaluation losses with increasing model size, validating the parameter efficiency of the generative approach. To further enhance performance, we propose a lightweight yet effective post-training pipeline that (1) enforces full-process SID alignment and (2) applies reinforcement learning with constrained decoding and hybrid rewards. Together, these techniques yield significant improvements in both ranking accuracy and candidate diversity.
% Our codes are available at 

\faGithub\ \url{https://github.com/AkaliKong/MiniOneRec}

\hspace{-0.3em} \raisebox{-0.2em}{
\includegraphics[height=0.9em]{figures/hf-logo.png}}
\url{https://huggingface.co/kkknight/MiniOneRec}


\end{boxedabstract}

%\begin{center}
%  \includegraphics[width=0.66\textwidth]{figures/scaling.png}
%  \captionsetup{type=figure}       
%  \captionof{figure}{Scaling curves from 0.5B to 7B parameters.}
%  \label{fig:scaling}
%\end{center}
% \end{abstract}

\begin{center}
\begin{minipage}[t]{0.46\textwidth}
\centering
\includegraphics[width=\linewidth]{figures/scaling_big.png}
\end{minipage}\hfill
\begin{minipage}[t]{0.46\textwidth}
\centering
\includegraphics[width=\linewidth]{figures/llmmatter_big.png}
\end{minipage}
\captionsetup{type=figure}
\captionof{figure}{Left: Scaling curves from 0.5B to 7B parameters. Right: Effect of world knowledge on model performance: MiniOneRec-W/O ALIGN uses pretrained LLM weights but omits SID–text alignment, while MiniOneRec-Scratch is trained from random initialization and omits alignment.}
\label{fig:scaling}
\end{center}

\clearpage
\tableofcontents
% \thispagestyle{plain} % 去掉页眉（可选）
\clearpage

\section{Introduction}
The scaling behavior of recommendation models has recently attracted significant attention \citep{OneRec, mtgr, HSTU, Wukong, URM, RankMixer, HLLM, OnePiece}, driven largely by the success of large language models (LLMs) in demonstrating predictable performance gains with increased model size \citep{GPT3, gpt4, deepseek, qwen2.5, mistral-7b, llama3, Llama2, Deepseek-r1}.
However, traditional recommendation models have struggled to exhibit similar scaling laws \citep{SASRec, GRU4Rec, Caser, SGL, survey_seq}.
These systems typically allocate the majority of parameters to large embedding tables for storing user and item representations, while using only inner-product or shallow scoring networks for final predictions.
Such an embedding-heavy design leads to performance plateaus: even as embedding dimensions or table sizes increase, improvements diminish quickly beyond moderate scales \citep{DIN, DCN, DCNv2, DeepFM}.

Generative recommendation offers a fundamental paradigm shift.
By compressing items into sequences of discrete Semantic IDs (SIDs) through quantization techniques \citep{rqkmeans, rqvae}, it uses compact vocabularies and redirects the bulk of parameters toward deep autoregressive Transformers that generate SID sequences \citep{TIGER, LCRec, TOKENRec, OneRec, OnePiece}.
This design enables scaling behaviors more characteristic of language models, where increased depth and capacity translate to consistent performance gains \citep{OneRec}.

Recent industrial deployments have demonstrated the promise of this generative paradigm.
OneRec \citep{OneRec, onerec-tech-v1} achieves significant improvements on Kuaishou's platform with 400 million daily active users, while OneRec-V2 \citep{onerec-tech-v2} further advances through lazy decoder architecture and preference alignment. 
OnePiece \citep{OnePiece} integrates LLM-style context engineering and reasoning into both retrieval and ranking models of industrial cascaded pipelines.
However, these results rely on massive proprietary datasets and remain closed-source, leaving critical questions unanswered for the research community:
\begin{itemize}
    \item Do the scaling advantages of generative recommendation transfer to public datasets?
    \item What is the minimal post-training recipe needed to achieve strong performance?
\end{itemize}



In this work, we present \textbf{MiniOneRec}, the first fully open-source framework for generative recommendation, which provides an end-to-end workflow encompassing SID generation, model supervised fine-tuning (SFT), and recommendation-driven reinforcement learning (RL). The release includes complete source code, reproducible training pipelines, and publicly available model checkpoints.
We implement SID construction using residual quantized variational autoencoder (RQ-VAE) \citep{rqvae}, in line with prior work such as TIGER \citep{TIGER, OneRec}.
We post-train Qwen-based (backbone) generative models at multiple scales, ranging from 0.5B to 7B parameters, on public benchmarks from the Amazon Review Dataset.
Following OneRec's successful two-stage post-training pipeline, we adopt SFT on user-item interaction sequences, followed by group relative policy gradient (GRPO) \citep{Deepseek-r1, deepseekmath} for alignment.
The main contributions of our work are summarized as follows:
\begin{itemize}
    \item \textbf{Validating the Scaling Laws of Generative Recommendation}:
    We present the first systematic investigation of how generative recommendation models scale on public data, to the best of our knowledge. Specifically, we train four MiniOneRec variants, spanning 0.5\,B to 7\,B parameters, on the Amazon Review corpus~\citep{amazon} with identical data splits and hyper-parameter settings. Both the final training loss and the held-out evaluation loss consistently decrease as model size grows (see Figures \ref{fig:scaling} Left and \ref{fig:scaling_eval_loss}), highlighting the superior parameter efficiency of the generative paradigm.
    \item \textbf{Optimizing Post-training Strategies for Generative Recommendation}: 
    We design a lightweight yet comprehensive post-training pipeline that tightly couples full-process SID alignment with reinforced preference optimization. First, we verify the influence of world knowledge on the model’s generative recommendation performance (see Figure \ref{fig:scaling} Right) and augment the vocabulary with dedicated SID tokens and enforce auxiliary alignment objectives throughout the two-stage optimization process. Second, during the RL phase, we shape the generation process itself: invalid tokens are masked to guarantee that every step produces a legal item, beam search is employed for efficient exploration of diverse candidates, and the reward signal combines rule-based accuracy with a ranking-aware penalty that pushes the model away from hard negatives. As a result, MiniOneRec offers a concise yet effective recipe for bringing RL with value-based ranking into generative recommendation, simultaneously improving candidate diversity and ranking fidelity.
\end{itemize}

The remainder of this report is organized as follows:
Section \ref{bg} reviews background on generative recommendation and the application of LLMs with RL in recommender systems.
Section \ref{method} presents our proposed MiniOneRec, including task formulation, item tokenization, and post-training pipeline.
Section \ref{training_detail} describes the training details while Section \ref{results} provides comprehensive experimental results.
Section \ref{conclusion} concludes with discussions of future directions.

% \clearpage
\tableofcontents
% \thispagestyle{plain} % 去掉页眉（可选）
\clearpage

\section{Introduction}
The scaling behavior of recommendation models has recently attracted significant attention \citep{OneRec, mtgr, HSTU, Wukong, URM, RankMixer, HLLM, OnePiece}, driven largely by the success of large language models (LLMs) in demonstrating predictable performance gains with increased model size \citep{GPT3, gpt4, deepseek, qwen2.5, mistral-7b, llama3, Llama2, Deepseek-r1}.
However, traditional recommendation models have struggled to exhibit similar scaling laws \citep{SASRec, GRU4Rec, Caser, SGL, survey_seq}.
These systems typically allocate the majority of parameters to large embedding tables for storing user and item representations, while using only inner-product or shallow scoring networks for final predictions.
Such an embedding-heavy design leads to performance plateaus: even as embedding dimensions or table sizes increase, improvements diminish quickly beyond moderate scales \citep{DIN, DCN, DCNv2, DeepFM}.
Generative recommendation offers a fundamental paradigm shift.
By compressing items into sequences of discrete Semantic IDs (SIDs) through quantization techniques \citep{rqkmeans, rqvae}, it uses compact vocabularies and redirects the bulk of parameters toward deep autoregressive Transformers that generate SID sequences \citep{TIGER, LCRec, TOKENRec, OneRec, OnePiece}.
This design enables scaling behaviors more characteristic of language models, where increased depth and capacity translate to consistent performance gains \citep{OneRec}.

Recent industrial deployments have demonstrated the promise of this generative paradigm.
OneRec \citep{OneRec, onerec-tech-v1} achieves significant improvements on Kuaishou's platform with 400 million daily active users, while OneRec-V2 \citep{onerec-tech-v2} further advances through lazy decoder architecture and preference alignment. 
\kxy{OnePiece \citep{OnePiece} integrates LLM-style context engineering and reasoning into both retrieval and ranking models of industrial cascaded pipelines}.
% \wjc{TODO: more prior work from industry.}
However, these results rely on massive proprietary datasets and remain closed-source, leaving critical questions unanswered for the research community:
\begin{itemize}
    \item Do the scaling advantages of generative recommendation transfer to public datasets?
    % \item How do generative models perform across different scales?
    \item What is the minimal post-training recipe needed to achieve strong performance?
\end{itemize}
% Without open implementations and released models, the community cannot regoriouly validate these claims or build upon them systematically.

% In this work, \textbf{we present MiniOneRec, the first fully open-source generative recommendation framework with complete code, training pipelines, and released model checkpoints}.
\kxy{In this work, \textbf{we present MiniOneRec, the first fully open-source generative recommendation framework that provides an end-to-end workflow encompassing SID generation, model Supervised Fine-Tuning (SFT), and recommendation-driven Reinforcement Learning.} The release includes complete source code, reproducible training pipelines, and publicly available model checkpoints.}
We implement SID construction using Residual Quantized Variational Autoencoder (RQ-VAE) \citep{rqvae}, in line with prior work such as TIGER \citep{TIGER, OneRec}.
We post-train Qwen-based (backbone) generative models at multiple scales, ranging from 0.5B to 7B parameters, on public benchmarks from the Amazon Review Dataset.
Following OneRec's successful two-stage post-training pipeline, we adopt SFT on user-item interaction sequences, followed by Group Relative Policy Gradient (GRPO) \citep{Deepseek-r1, deepseekmath} for alignment.
The main contributions of our work are summarized as follows:
\begin{itemize}
    \item \textbf{Validating the Scaling Laws of Generative Recommendation}:
    We present the first systematic investigation of how generative recommendation models scale on public data. Specifically, we train four MiniOneRec variants, spanning 0.5\,B to 7\,B parameters, on the Amazon Review corpus~\citep{amazon} with identical data splits and hyper-parameter settings. Both the final training loss and the held-out evaluation loss consistently decrease as model size grows (see Figure~\ref{fig:scaling};\ref{fig:scaling_eval_loss}), highlighting the superior parameter efficiency of the generative paradigm.
    \item \textbf{Optimizing Post-training Strategies for Generative Recommendation}: 
    We design a lightweight yet comprehensive post-training pipeline that tightly couples full-process SID alignment with reinforced preference optimization. First, we augment the vocabulary with dedicated SID tokens and enforce auxiliary alignment objectives throughout the two-stage optimization process. Second, during the reinforcement learning (RL) phase, we shape the generation process itself: invalid tokens are masked to guarantee that every step produces a legal item, beam search is employed for efficient exploration of diverse candidates, and the reward signal combines rule-based accuracy with a ranking-aware penalty that pushes the model away from hard negatives. The resulting framework, MiniOneRec, offers a concise yet effective recipe for bringing reinforcement learning with value-based ranking into generative recommendation, simultaneously improving candidate diversity and ranking fidelity.
\end{itemize}

The remainder of this report is organized as follows:
Section \ref{bg} reviews background on generative recommendation and the application of LLMs with RL in recommendation systems.
Section \ref{method} presents our proposed MiniOneRec, including task formulation, item tokenization, and post-training pipeline.
Section \ref{training_detail} describes the training details while Section \ref{results} provides comprehensive experimental results.
Section \ref{conclusion} concludes with discussions of future directions.
% % \section{Introduction}
\clearpage

% Recommender systems aim to recover users’ latent interests from their past interactions. Inspired by the rapid progress of generative models—especially large language models (LLMs) \cite{GPT3,gpt4,qwen2.5,deepseek}—the community has begun to recast recommendation as a generative task. Under this paradigm, known as generative recommendation \cite{TIGER,LCRec,TOKENRec,OneRec,mtgr}, the model directly produces a sequence of item identifiers in an autoregressive fashion, replacing the conventional score-then-rank pipeline. 
% Representing items with Semantic IDs (SIDs) lets the generative recommender share parameters across semantically related items, cover billion-scale catalogs with a compact vocabulary, and naturally align with the hierarchical reasoning style of LLMs.

The powerful sequence modeling capabilities of large language models (LLMs) have enabled their adaptation to recommender systems \citep{BIGRec, AlphaRec, AlphaFuse, LLM2REC, survey_LLM_Rec, survey_seq}, with generative recommendation emerging as a promising direction that leverages autoregressive generation for item prediction \citep{TIGER, LCRec, TOKENRec, HSTU, OneRec, mtgr}.
% This paradigm centers on generating sequences of Semantic IDs (SIDs)---discrete tokens that encode item semantics through quantization of continuous embeddings \citep{rqvae, rqkmeans}.
% By promoting token sharing across semantically related items, SIDs facilitate efficient handling of large-scale catalogs while naturally aligning with the step-by-step (chain-of-thought) reasoning paradigm of LLMs \citep{TIGER, rqvae, rqkmeans, OneRec, bettersid}.
\kxy{Notably, unlike Large Embedding Models~(LEMs) that devote the majority of their parameters to embedding tables, OneRec\citep{OneRec} shows that increasing network depth is an effective way to learn Semantic IDs~(SIDs)—discrete tokens that represent item semantics through the quantization of continuous embeddings~\citep{rqvae,rqkmeans}. To a certain extent, this finding corroborates the emerging scaling law in recommendation: deeper networks, while keeping the vocabulary size compact, can absorb richer semantic signals and consequently yield noticeably better recommendation performance.}


Upon scrutinizing prior studies on generative recommenders, we can summarize a common \za{training} pipeline:
(1) \za{Beginning with items'} textual descriptions or embeddings, a quantization method transforms continuous vectors into SIDs.
(2) The model is then trained to generate these SIDs in an end-to-end \za{manner, typically following} two primary approaches: \textit{scratch-trained recommenders} that train Transformer \citep{Transformer} from scratch on user interaction sequences \citep{TIGER, HSTU, OneRec, mtgr}, and \textit{SID-aligned Recommenders} that adapt pretrained LLMs \za{from the language space into SID space} through supervised fine-tuning (SFT) \citep{LCRec, TOKENRec}.
\za{Although these methods achieve promising performance, the integration of reinforcement learning (RL) for deeper alignment with user interaction sequences remains relatively underexplored \citep{sdpo}.} 

\za{
Inspired by the recent success of the \textbf{SFT-then-RL} paradigm \citep{deepseekmath, SFT-then-RL, o1}, we aim to establish a generative recommendation framework that adapts this paradigm to the unique characteristics of recommendation tasks. 
Specifically, we first explore applying Group Relative Policy Optimization (GRPO) \citep{Deepseek-r1} after the initial SFT phase, considering both scratch-trained and SID-aligned recommenders.
Our preliminary experiments (\cf Figure \ref{fig:sft_vs_rl}) demonstrate the effectiveness of the SFT-then-GRPO paradigm in generative recommendation, consistently outperforming the SFT-only baselines.
Building on this, we further investigate two critical limitations that arise when directly transferring the standard SFT-then-RL paradigm to recommendation, which constrain its full potential in this domain:
}

% Recently, the \textbf{SFT-then-RL} post-training paradigm has shown promise for generative models \citep{deepseekmath, Deepseek-r1, SFT-then-RL, o1}.
% Following this trend, we first explore applying this approach to generative recommenders by incorporating Group Relative Policy Optimization (GRPO) \citep{Deepseek-r1} after the initial SFT phase.
% Our preliminary experiments (\cf Figure \ref{fig:sft_vs_rl}) confirm the effectiveness of this paradigm, showing consistent improvements over SFT-only baselines.

% However, upon deeper investigation, we identify two critical limitations in directly adapting the standard SFT-then-RL approach that constrain its full potential:

% ======================


% Upon scrutinizing previous studies on generative recommenders, we can summarize a common pipeline:
% (1) Starting from the text of an item—or an embedding from another modality—then apply a quantization method to turn the continuous vector into an ordered list of code words, called SIDs.
% (2) A model is then trained to generate these semantic IDs in an end-to-end fashion. There are two major variants:
% \begin{enumerate}[label=(\alph*),leftmargin=1.5em]
% \item \textbf{Scratch-trained Recommender:} a Transformer is pre-trained from scratch on user-interaction sequences composed of SIDs \citep{TIGER, OneRec, mtgr}.
% \item \textbf{SID-aligned Recommender:} a pre-trained LLM is aligned with the item representation space through supervised fine-tuning (SFT) \citep{LCRec, TOKENRec}.
% \end{enumerate}
% (3) At inference time, the model uses beam search together with constrained decoding to generate valid SID sequences.

% Recently, the \textbf{SFT-then-RL} paradigm has proven highly effective for generative models \citep{deepseekmath, Deepseek-r1, SFT-then-RL, o1}. In this work, we first follow the LLM-based training pipeline: we perform SFT to align LLM with the item representation space, and then apply the Group relative policy optimization (GRPO) algorithm to raise the performance ceiling further. 
% As illustrated in Figure \ref{fig:sft_vs_rl}, the SFT-then-RL paradigm has surpassed the earlier SFT-only approach in generative recommendation. 
% Although this two-stage strategy represents meaningful progress, it still suffers from several limitations:

 \begin{figure*}[t]
  \centering
  \begin{subfigure}[t]{0.34\textwidth}
    \centering
    \includegraphics[width=\textwidth]{figures/sft_vs_sft_rl.png}
    \caption{Performance under Different Training Paradigms.}
    \label{fig:sft_vs_rl}
\end{subfigure}\hspace{2mm}
\begin{subfigure}[t]{0.43\textwidth}
    \centering
    \raisebox{15pt}{% 
        \includegraphics[width=\textwidth]{figures/pie.png}}
    \caption{Sample hit reward distribution.}
    \label{fig:pie}
\end{subfigure}
  \caption{Results on the Industrial dataset. Figure \ref{fig:sft_vs_rl} illustrates the performance, with respect to HitRatio@10 under different training paradigms. Figure \ref{fig:pie} shows the proportion of RL-sampled outputs from the SFT-initialized model that receive a non-zero reward.}
  \label{fig:teaser}
    \vspace{-10pt}
\end{figure*}


\begin{itemize}[leftmargin=*]
    % \item \textbf{Superficial SID understanding via world knowledge.}
    \label{llmmatter}
    \item \textbf{Limited SID Understanding in Alignment}.
    The SFT alignment process \za{constrains} LLMs to project outputs \za{into a closed SID vocabulary} through supervised training on user-item interaction sequences. 
    Such rigid constraint \za{mainly encourages superficial pattern matching rather than true semantic understanding of SIDs.}
    As illustrated in Figure \ref{fig:superficial_align}, even after \za{SFT the model often fails to exploit SID histories correctly: it either resorts to generic responses citing insufficient user information, or produces} lengthy, repetitive item descriptions \za{indicative of collapsed generation.}
    %when presented with SID-formatted interaction histories, models frequently fail to leverage these representations meaningfully. Instead, they either produce generic responses claiming insufficient user information, or generate lengthy, repetitive item descriptions that suggest a collapse in their generative capabilities.
    While the world knowledge in pretrained LLMs should be beneficial for understanding item semantics and user preference (\cf Figure \ref{fig:sft_vs_rl}), the current alignment pipeline exploits it only \za{superficially}, highlighting the need for a more systematic approach that integrates SID understanding throughout the entire training process.

    % \item \textbf{Ineffective Reward Assignment in RL}.
    % The standard rule-based RL training treats each SID sequence as an indivisible unit---if any token in a generated sequence fails to match the ground-truth, the entire sequence \za{is penalized with} zero reward. 
    % This binary reward mechanism overlooks the rich relational structure among SIDs, unable to distinguish between near-correct and completely irrelevant predictions.
    % The resulting sparse rewards (\cf Figure \ref{fig:pie}) \za{deprives the model of useful learning signals on hard cases}---those rollouts that \za{narrowly} miss exact matches and thus receive zero advantages \citep{DAPO}---which represent the most critical cases for developing genuine reasoning capabilities over the SID taxonomy.
    % Consequently, RL optimization \za{tends to} reinforce patterns already memorized during pre-training and SFT \citep{sftmemorize,RLReallyLearn,RLcriticR1}, rather than learning to navigate semantic relationships that would enable better performance on challenging cases.
    \kxy{
    \item \textbf{Unique generation space}.
    In generative recommendation, the valid item space is much narrower than in open-ended language generation. If decoding is not explicitly constrained, the model frequently produces invalid items that are nonexistent in the item corpus \citep{D3}. Furthermore, the constrained generation space leads to a steep token probability distribution, which makes the model prone to repeatedly generating the same items across multiple samples. As a result of both invalid and repetitive generations, the model encounters difficulties in both inference and sampling: standard sampling strategies frequently yield redundant items, substantially reducing sampling efficiency during training. Consequently, at each training step, the model is exposed to only a very limited range of negative samples, which restricts the diversity of ranking signals and results in inefficient optimization.

    \item  \textbf{Sparse ranking supervision}.
    Rule-based reward modeling in reinforcement learning with verifiable rewards (RLVR) \citep{Deepseek-r1} functions as a binary correctness signal, assigning a reward of one to the correct item and zero to all others. In recommendation, however, user interactions are inherently sparse, and the vast majority of items in the corpus remain unobserved (Wei et al., 2024). As a result, nearly all sampled items are uniformly treated as negatives, each assigned an identical reward of zero. This leads to sampling groups where only the single target item is distinguished, while all other items collapse into the same reward value. Such sparsity provides only weak supervision signals during optimization and leaves the model without fine-grained ranking feedback.
    }
    
\end{itemize}



% \begin{itemize}[leftmargin=*]
% %     \item \textbf{Superficial SID understanding via world knowledge.} 
% %     To quantify the benefit of world knowledge, we conduct a controlled study on the large-scale Amazon-Book dataset \citep{amazon}, contrasting:  
% % (1) a Scratch-trained Recommender and  
% % (2) a SID-aligned Recommender.
% % As shown in Figure \ref{fig:sft_vs_rl}, the latter consistently surpasses the former on multiple evaluation metrics, indicating that the factual and commonsense knowledge in LLMs is beneficial for generative recommendation.  \label{pretrainwin}
% % Prior work has further demonstrated that explicitly aligning the textual embedding space with a SID space can deliver additional gains for SID-aligned recommenders \citep{LCRec,TOKENRec}.   
% % Unfortunately, the standard SFT alignment pipeline only forces the LLM to project its outputs onto a closed set of SIDs. 
% % Such a rigid constraint encourages the model to memorize only surface patterns and results in catastrophic forgetting of its general capabilities \citep{sftruin,sftmemorize, RLvsSFT}, as illustrated in Figure \ref{fig:superficial_align}. Consequently, the world knowledge learned during pre-training is exploited only in a superficial manner, leaving the final recommendation performance sub-optimal.
%     \item \textbf{Coarse-grained use of SID information.} In the prevailing rule-based RL, SID prediction is treated as an exact-match classification task: a positive reward is given only when the SID produced by the model is identical to the ground-truth SID, while all other SIDs—regardless of their semantic proximity—receive the same penalty.  
% This binary feedback overlooks the rich relational structure among SIDs and produces an coarse reward landscape. As illustrated in Figure \ref{fig:pie}, the SFT-initialized model still receives sparse rewards during RL. For hard instances, positive rewards become almost unattainable, making effective credit assignment nearly impossible. Consequently, the apparent performance gains of RL mostly come from polishing patterns already memorised during pre-training and SFT, rather than from truly leveraging the SID taxonomy to raise the capability ceiling~\citep{sftmemorize,RLReallyLearn,RLcriticR1}.

% \end{itemize}
To address the above limitations, we introduce \textbf{MiniOneRec}, a generative recommendation framework that
enhances both SID comprehension and reward utilization throughout the \za{SFT-then-RL} process.
% not only tightly integrates SIDs with the world knowledge encoded in pretrained LLMs, but also exploits the hierarchical structure of SIDs at a finer granularity to further boost recommendation quality.
MiniOneRec is built on two key components:  
\begin{itemize}[leftmargin=*]
    \item \textbf{Full-Process SID Alignment.} 
    We embed a set of alignment objectives throughout the entire SFT-then-RL pipeline to achieve deeper SID alignment. 
    Specifically, we integrate discrete SID tokens into the LLM's vocabulary and introduce a series of auxiliary alignment tasks (\eg explicit item title to SID mapping) that are \za{enforced during both SFT and RL phases.} 
    \za{This ensures the model fully internalizes the structural semantics of SIDs.}
    % applied consistently across both SFT and RL phases to ensure the model fully internalizes the structural semantics of SIDs.
    % Specifically, we integrate the discrete SID tokens into the LLM’s vocabulary. In addition to the standard semantic tasks, we devise a series of alignment tasks (\emph{e.g.} explicitly map each item title to its SID) and
    % apply them across both SFT and RL phases
    % apply them throughout the entire SFT-then-RL pipeline to deepen the model’s understanding of SIDs.
    % The treatment enables the model to internalise the structural semantics of SIDs, translating into measurable gains on downstream generative-recommendation benchmarks.
    % \item \textbf{SID-Navigated Reinforcement Learning.}
    % We develop a novel RL framework that leverages SID hierarchy to provide more informative learning signals:
    % \begin{enumerate}[label=(\alph*),leftmargin=1.5em]
        %\item \textbf{SID-Prefix Curriculum Sampling.} Let the ground-truth item be a hierarchical token sequence $e^{pos}=(s^\star_a,s^\star_b,s^\star_c)$. Inspired by curriculum learning, we define a scheduling function \(f(t,e^{pos})\) that truncates 
% $e^{pos}$ and progressively shortens the retained prefix as the training step $t$ increases. 
% The resulting truncated prefix is concatenated with the original input to construct a new prompt.
%The policy then produces a continuation and the reward is computed accordingly.  
%In this way, the model starts by predicting lower-level tokens conditioned on higher-level prefixes and gradually transitions to prefix-free generation, thereby achieving autonomous exploration and performance gains on difficult samples.
        
        % \item \textbf{SID-Level Reward Modeling.}
        %We employ GRPO to optimize the LLM with fine-grained, SID-level rewards. For hard samples whose original rewards are often zero, we treat the step-by-step generation of the SID path $\langle a \rangle \!\to\! \langle b \rangle \!\to\! \langle c \rangle$ as an intermediate reasoning process. 
        % \za{Suppose the ground-truth SID token be $e^{pos}=(s^\star_a,s^\star_b,s^\star_c)$ and the model output be $e=(s_a,s_b,s_c)$}.
        %A partial reward is granted according to the deepest level that still matches the ground truth, \za{\ie $\kappa(e,{e}^{pos})=\max\{\,k:\, e_{a:k}=e^\star_{a:k}\,\}$, preventing zero-reward collapse on hard examples.}
        %\end{enumerate}

    \kxy{\item \textbf{Reinforced Preference Optimization}.
    For generation, we adopt constrained decoding by masking invalid tokens at each step, ensuring that only valid items are produced and fundamentally distinguishing the output space from that of open-ended language tasks.For sampling, we employs beam search to efficiently generate diverse candidate items in a single pass, guaranteeing both sampling efficiency and exposure to informative negatives. For reward modeling, we augment rule-based accuracy rewards with ranking rewards, which assign additional penalties to hard negatives according to their generation probabilities. Altogether, MiniOneRec offers a principled yet practical way to bring RLVR into generative recommendation, enhancing both the quality of negatives and the fidelity of ranking signals.}
\end{itemize}

\begin{figure}
    \centering
    \includegraphics[width=1.0\textwidth]{figures/limited_understanding.png}
    \vspace{-10pt}
    \caption{Only-SFT-Align case study.  
The model is fed a SID-formatted interaction history and asked to describe the next item. The Only-SFT-Align LLM fails to interpret the SID tokens (Answer 1) or generates verbose, repetitive, and disordered text (Answer 2), underscoring the shallow alignment achieved by SFT alone. By contrast, MiniOneRec correctly understands the SID sequence and produces a concise, coherent item description.}
    \label{fig:superficial_align}
    \vspace{-10pt}
\end{figure}

We evaluate MiniOneRec on two public benchmark datasets \citep{amazon} and compare it with a wide range of leading traditional sequential recommenders, generative recommenders, and several recent LLM-based baselines. The results show that MiniOneRec consistently outperforms all competitors on standard recommendation metrics, confirming its effectiveness.
% \emph{A detailed survey of related work is deferred to Appendix~\ref{appendix:related_work}.}
\section{Background and Related Work}
\label{bg}
\subsection{Generative Recommendation}

In the last few years, formulating recommendation as a sequence generation problem has become a vibrant research direction \citep{TIGER, GenRec-Direction}.  
Under this view, a recommender is trained to predict the tokens of the next item in an end-to-end manner, typically with a Transformer backbone \citep{Transformer}.  
Such a design removes the rigid \slh{multi-stage} “matching–ranking” pipeline in traditional retrieval-based recommender systems \citep{SASRec, GRU4Rec} and therefore pushes the performance upper bound.

Early explorations illustrate two key components: (1) turning an item into a discrete code (SIDs) and (2) letting the model produce the code.  
TIGER~\citep{TIGER} adopts RQ-VAE~\citep{rqvae} to map textual embeddings of titles and descriptions into SIDs.  
HSTU~\citep{HSTU} introduces a streaming architecture that is friendly to high-cardinality and non-stationary logs.  
LC-Rec~\citep{LCRec} aligns an LLM with the SIDs through multi-task learning so that the model can understand these symbols during generation.  

Subsequent work focuses on designing better codes.  
RecForest~\citep{recforest} clusters items via hierarchical $k$-means and uses the cluster indices as tokens.  
EAGER~\citep{eagar} and TokenRec~\citep{TOKENRec} fuse collaborative and semantic evidence directly into the tokenizer.  

At industrial scale, generative recommenders have started to replace heavy cascade systems.  
MTGR~\citep{mtgr} keeps the original DLRM features, adds user-level compression, and accelerates both training and inference.  
OneRec~\citep{onerec-tech-v2} reduces serving cost through a lazy decoder-only layout and stabilizes optimization with an improved RL algorithm.  
OnePiece~\citep{OnePiece} discovers that performing inference in the latent space can further improve generative recommendation performance.

\subsection{LLM and RL}

RL aims at maximizing cumulative reward through repeated interaction~\citep{RLsurvey}.  
When fine-tuning LLMs, RL with Human Feedback (RLHF) has become a standard recipe; PPO~\citep{PPO} is the most common optimizer but is memory-intensive for billion-scale parameters.  
To reduce cost, Direct Preference Optimization (DPO)~\citep{DPO} removes the separate value network and maximizes the log-likelihood gap between preferred and dispreferred outputs.  
S-DPO~\citep{sdpo} adapts this idea to recommendation by treating softmax-based negative sampling as implicit pairwise preference.  
Nevertheless, preference-based objectives are off-policy and may converge prematurely.

Light-weight on-line methods have therefore been proposed.  
GRPO~\citep{deepseekmath} normalizes rewards within a small batch of roll-outs and replaces the learned reward model by rule-based signals, achieving strong gains in maths and code generation.  

% \section{Methods}
\slh{
\subsection{Item Tokenization}

Following the standard item tokenization pipeline in SID-based generative recommender systems, such as TIGER \cite{TIGER} and OneRec \cite{OneRec}, we adopt RQ-VAE \cite{rqvae} as the item tokenization method. }
For each item $i$, we concatenate its title and textual description and feed the resulting sentence into a frozen content encoder to produce a $d$-dimensional semantic vector $\mathbf{x}\in\mathbb{R}^{d}$.
% $\mathbf{x}\in\mathbb{R}^{d}$.
The continuous vector is then discretized with the RQ-Kmeans algorithm~\citep{rqkmeans}, which builds a hierarchy of codebooks by recursively clustering the residuals.

Let $\tilde{\mathbf{M}}=[\mathbf{x}_1;\ldots;\mathbf{x}_N]\in\mathbb{R}^{N\times d}$ be the matrix that stacks the embeddings of all $N$ items.  
We initialize $\mathbf{R}^{(1)}=\tilde{\mathbf{M}}$. For each layer $l\in\{1,\ldots,L\}$, where $L$ is the number of hierarchical levels in the semantic codebook, we learn a codebook
$\mathbf{C}^{(l)}=\{\mathbf{c}^{(l)}_{k}\}_{k=1}^{K_l}$ by running K-means with $K_l$ centroids on the current residuals $\mathbf{R}^{(l)}$:
\[
\mathbf{C}^{(l)}=\text{K-means}\bigl(\mathbf{R}^{(l)},K_l\bigr).
\]
%
For item $i\;(1\le i\le N)$, the index of the nearest centroid is obtained via
\[
s^{(l)}_{i}=\arg\min_{k}\lVert\mathbf{R}^{(l)}_{i}-\mathbf{c}^{(l)}_{k}\rVert_{2},
\]
where $\lVert \cdot \lVert$  denotes the Euclidean norm. The residual is updated as
\[
\mathbf{R}^{(l+1)}_{i}= \mathbf{R}^{(l)}_{i}-\mathbf{c}^{(l)}_{s^{(l)}_{i}}.
\]
After $L=3$ layers we obtain a coarse-to-fine set of semantic identifiers,
$
\{s^{(1)}_{i},s^{(2)}_{i},s^{(3)}_{i}\},
$
which serves as the unique token sequence for item $i$ and will be consumed by the recommender for progressive generation.

% \begin{equation}
    %L_{RECON} = ||e-\hat{e}||_{2}^{2}
% \end{equation}

% \begin{equation}
% \mathcal{L}_{\mathrm{RQ}}=\sum_{i=0}^{m-1}\left\|\operatorname{sg}\left[\boldsymbol{r}_{i}\right]-\boldsymbol{e}_{c_{i}}^{i}\right\|_{2}^{2}+\beta\left\|\boldsymbol{r}_{i}-\operatorname{sg}\left[\boldsymbol{e}_{c_{i}}^{i}\right]\right\|_{2}^{2},
% \end{equation}

% \begin{equation}
% L_{rq-vae} = L_{RECON} + L_{RQ}
% \end{equation}


\subsection{Model Optimization}

We fine-tune our policy with the GRPO algorithm \citep{deepseekmath,deepseek}, which leverages the relative quality of multiple responses generated for the same prompt. Concretely, for each input $x\!\sim\!D$, we roll out the current policy $\pi_{\theta_{\text{old}}}$ $G$ times to obtain a set of candidates $\mathcal{Y}(x)=\{y^{(1)},\dots,y^{(G)}\}$.  Each candidate $y^{(i)}$ receives a scalar reward $R_i$, and the advantage is computed by normalizing the rewards \emph{within the group}
\begin{equation}
    \hat{A}_i=\frac{R_i-\operatorname{mean}(R_{1:G})}{\operatorname{std}(R_{1:G})},
    \label{eq:grpo_adv}
\end{equation}
where $R_{1:G}=\{R_1,\dots,R_G\}$. This groupwise normalization recenters advantages at zero and rescales them to unit variance, thereby turning each prompt into a self-contained comparison game and reducing gradient variance.  The new policy $\pi_\theta$ maximizes the clipped surrogate
\begin{equation}
\begin{aligned}
J_{\text{GRPO}}(\theta)=\;
&\mathbb{E}_{x\sim D,\; y^{(i)}\sim\pi_{\theta_{\text{old}}}}
\Biggl[
\frac{1}{G}\sum_{i=1}^{G}\frac{1}{|y^{(i)}|}
\sum_{t=1}^{|y^{(i)}|}
\Bigl\{
\min\!\bigl(r_{i,t}\hat{A}_{i,t},\,
\operatorname{clip}(r_{i,t},1-\epsilon,1+\epsilon)\hat{A}_{i,t}\bigr)
\\[4pt]
&\qquad
-\beta\,\mathrm{KL}\!\bigl[\pi_\theta\|\pi_{\text{ref}}\bigr]
\Bigr\}
\Biggr],
\end{aligned}
\label{eq:grpo_obj}
\end{equation}
where 
$r_{i,t}=
\tfrac{\pi_\theta(y^{(i)}_t\,|\,x,y^{(i)}_{<t})}
      {\pi_{\theta_{\text{old}}}(y^{(i)}_t\,|\,x,y^{(i)}_{<t})}$
is the per-token importance ratio, $\epsilon$ is the clipping threshold, and $\beta$ balances task reward against a KL penalty, keeping the updated policy close to a reference model $\pi_{\text{ref}}$.


\subsection{Task formulation}
Generative recommendation reformulates the recommendation problem as a sequence generation task.
% where an autoregressive sequence model is trained to output a series of item identifiers that constitute personalized recommendations.
Let $H_u$ denote the interaction history of user $u$, sorted in chronological order.  
Each item $i \in H_u$ is represented by a $3$–level SID tuple $\{s^{(1)}_{i},\,s^{(2)}_{i},\,s^{(3)}_{i}\}$.  
Given $H_u$, the generative recommender $\pi_{\theta}$, parameterized by $\theta$, is trained to predict an item $i^{\text{pos}}$ that best matches the preferences of user $u$ from the item set.  
During inference, we employ beam search to generate the top-$k$ candidates and evaluate the model with standard generative-recommendation metrics.
\section{Modeling}
\label{method}
In this section, we present the modeling strategies of \textbf{MiniOneRec}, as illustrated in Figure \ref{fig:framework}. 
MiniOneRec first converts the textual information into SIDs with the RQ-VAE tokenizer (Section \ref{sec:tokenization}). 
For better incorporating the huge world knowledge within LLMs \citep{LlaRA}, MiniOneRec further introduces the alignment with LLMs as one expansion of the original OneRec architecture (Section \ref{sec:alignment_llm}). 
During the training stage, MiniOneRec is optimized with next token prediction first, and then followed by reinforeced preference optimization (Section \ref{sec:optimization}). 


\begin{figure}
    \centering
    \includegraphics[width=0.85\textwidth]{figures/MiniOneRec_framework.pdf}
    \vspace{-10pt}
    \caption{MiniOneRec framework. RQ-VAE builds the item SID codebook. We then perform SFT to warm up the LLM and obtain an initial alignment. In RL, beam search with constrained decoding, thereby the model sequentially produces a ranked list of distinct, valid SIDs. GRPO updates the policy, and SID alignment is enforced end-to-end. This alignment objective is preserved throughout both the SFT and RL stages, fostering deeper semantic understanding.}
    \label{fig:framework}
    \vspace{-10pt}
\end{figure}

% To address the key shortcomings of the existing SFT-then-RL pipeline for generative recommendation, we propose \textbf{MiniOneRec}.  
% The framework elevates the performance ceiling by aligning the LLM with the SID space throughout the entire training process and applying preference optimization in the RL stage, as illustrated in Figure \ref{fig:framework}.

\subsection{Task formulation}

We first formulate recommendation as a sequence–generation problem.  
For every user $u$, the items he or she has interacted with are sorted chronologically to form a sequence
$H_u=\bigl[i_1,i_2,\dots,i_T\bigr]$.  
Each item $i_t$ in the sequence is encoded by a three–level structural ID, $\bigl\{c_0^{i_t},\,c_1^{i_t},\,c_2^{i_t}\bigr\}$. Such structural IDs are typically called SIDs, which preserve hierarchical semantics through quantization techniques with semantic embeddings \citep{TIGER}.
% since they are 
% so that hierarchical semantics are preserved.

A generative policy $\pi_\theta$, implemented as an autoregressive model with parameters $\theta$, reads the entire history $H_u$ and is trained to predict the next item $i^{+}$ that best matches the taste of user $u$ among all candidates in the catalog.  
During inference, the model recursively produces item tokens; we keep the $k$ most promising beams by the standard beam–search algorithm and return them as the recommendation list.  
Model performance is reported with the evaluation measures commonly adopted in generative recommendation.


\subsection{Item Tokenization} \label{sec:tokenization}

In SID-style generative recommenders, the first task is to convert each item into a sequence of discrete tokens. Following the practices of TIGER \citep{TIGER}, we employ RQ-VAE~\citep{rqvae} for this purpose. % The procedure is as follows: (1) For every item $i$, we concatenate its title and textual description to form a single sentence; (2) This sentence is passed through a frozen text encoder \citep{AlphaRec}, producing a $d$-dimensional semantic vector $\mathbf{x}\in\mathbb{R}^{d}$; (3) RQ-VAE then applies multi-stage residual quantization to $\mathbf{x}$, gradually approximating it and finally yielding a list of discrete codewords. These codewords are used as the item representation in the subsequent generative recommendation module.
Concretely, the pipeline is:

\begin{enumerate}
    \item For every item $i$, we concatenate its title and textual description to form a single sentence;  
    \item This sentence is passed through a frozen text encoder \citep{AlphaRec}, producing a $d$-dimensional semantic vector $\mathbf{x}\in\mathbb{R}^{d}$;  
    \item Apply RQ-VAE to $\mathbf{x}$. At each level $l$ ($0\le l<L$) we have a separate codebook $\mathcal{C}_l=\{\mathbf{e}^{(l)}_{k}\}_{k=1}^{K}$, where $K$ is the codebook size. We set $L{=}3$ and $K{=}256$, so each item is represented by three bytes; this choice provides $2^{24}$ possible codes, which is sufficient for catalogs containing hundreds of millions of products while keeping the vocabulary small. The residual is initialized as $\mathbf{r}_{0}=\mathbf{x}$ and updated by
    \[
       c_l=\arg\min_{k}\left\|\mathbf{r}_l-\mathbf{e}^{(l)}_{k}\right\|_2,\qquad 
       \mathbf{r}_{l+1}=\mathbf{r}_l-\mathbf{e}^{(l)}_{c_{l}}.
    \]
    \item Collect the indices $(c_0,\dots,c_{L-1})$ as the discrete token sequence for item $i$; these indices constitute the item tokens consumed by the subsequent generative recommender.
\end{enumerate}

The quantized latent is reconstructed by
\[
    \mathbf{z}_{\text{q}}=\sum_{l=0}^{L-1}\mathbf{e}^{(l)}_{c_l},
\qquad
    \hat{\mathbf{x}}=D(\mathbf{z}_{\text{q}}),
\]
where $D(\cdot)$ is a decoder. The codebooks, the encoder and the decoder are trained jointly, while the text encoder remains frozen. The loss is the sum of a reconstruction term and an RQ regularizer:
\[
\mathcal{L}(\mathbf{x})=
\underbrace{\left\|\mathbf{x}-\hat{\mathbf{x}}\right\|_2^{2}}_{\mathcal{L}_{\text{RECO}}}
+\underbrace{\sum_{l=0}^{L-1}\!\bigl(
      \|\mathrm{sg}[\mathbf{r}_l]-\mathbf{e}^{(l)}_{c_l}\|_2^{2}
      +\beta\,\|\mathbf{r}_l-\mathrm{sg}[\mathbf{e}^{(l)}_{c_l}]\|_2^{2}\bigr)}_{\mathcal{L}_{\text{RQ}}},
\]
with $\mathrm{sg}[\cdot]$ the stop-gradient operator and $\beta$ a hyper-parameter controlling the commitment term.

To prevent codebook collapse, we follow the warm-start trick in prior work \citep{rqvae, LCRec} and initialize each codebook with $k$-means centroids computed on the first training batch.


\subsection{Alignment with LLMs} \label{sec:alignment_llm}
LLMs possess extensive understanding of the world and human behaviors 
\citep{AlphaRec}, which can serve as a supplement to the collaborative signals in recommender systems \citep{RLMRec}.
Existing work has shown that linking an LLM’s world knowledge to SID representations noticeably strengthens generative recommendation \citep{LlaRA, LCRec}. Therefore, rather than training on SIDs alone as in prior work \citep{TIGER, OneRec, mtgr}, we introduce several alignment objectives that tie the language space to SID signals. Two major groups of tasks are employed:  
\begin{itemize}
    \item \textbf{Recommendation Tasks}: The LLM receives a time-ordered history together with a clear instruction and is asked to predict the SID of the next item the user might engage with.  
    \item \textbf{Alignment Tasks}: A collection of bridging tasks enforces a two-way mapping between natural language and SID space, grounding the discrete codes in text while injecting linguistic knowledge into their embeddings.  
\end{itemize}
Tasks from both groups are optimized jointly throughout the SFT stage and the subsequent RL stage. During RL, we adopt constrained decoding so that the model can only produce tokens from a predefined list containing every item’s SID and its canonical title. This constraint guarantees valid outputs and enables straightforward, rule-based reward computation. Detailed examples of the prompts are provided in the Appendix \ref{appendix:align_prompt}.


\subsection{Reinforced Preference Optimization} \label{sec:optimization}
After SFT, we further polish the policy with GRPO~\citep{deepseekmath,deepseek}. GRPO differs from classic RLHF in that it draws multiple candidates per prompt and normalizes rewards within the group, which reduces gradient variance.

Steps:
(1) For every prompt $x\!\sim\!D$, the frozen policy $\pi_{\theta_{\text{old}}}$ is rolled out $G$ times, yielding $\mathcal{Y}(x)=\{y^{(1)},\ldots,y^{(G)}\}$.  
(2) Each candidate $y^{(i)}$ is assigned a scalar score $S_i$.  
(3) Advantages are standardized inside the group:
\[
\hat{A}_i=\frac{S_i-\mu_{1:G}}{\sigma_{1:G}},
\]
where $\mu_{1:G}$ and $\sigma_{1:G}$ denote the mean and standard deviation of the $G$ rewards.

The surrogate objective becomes
\[
\small
\begin{aligned}
J_{\text{GRPO}}(\theta)=
\mathbb{E}_{x,y^{(i)}}
\biggl[
\frac{1}{G}\sum_{i=1}^{G}\frac{1}{|y^{(i)}|}
\sum_{t=1}^{|y^{(i)}|}
\!\Bigl(
\min\!\bigl(w_{i,t}\hat{A}_{i,t},\,
\operatorname{clip}(w_{i,t},1-\epsilon,1+\epsilon)\hat{A}_{i,t}\bigr)
-\beta\,\mathrm{KL}\!\bigl[\pi_\theta\|\pi_{\text{ref}}\bigr]
\Bigr)
\biggr],
\end{aligned}
\]
with token-level importance ratio  
$w_{i,t}= \dfrac{\pi_\theta(y^{(i)}_t\!\mid\!x,y^{(i)}_{<t})}
                 {\pi_{\theta_{\text{old}}}(y^{(i)}_t\!\mid\!x,y^{(i)}_{<t})}$.
The parameter $\epsilon$ controls clipping, while $\beta$ keeps the updated policy near a reference model through a KL term.


Applying RL with verifiable rewards (RLVR) to recommendation brings two obstacles:

\begin{itemize}
    \item \textbf{Unique generation space}. The action space is a closed set of item SIDs, orders of magnitude smaller than natural-language vocabularies. Re-sampling tends to produce duplicates, wasting computation. We therefore mix dynamic sampling \citep{DAPO} with constrained beam search to enlarge coverage while keeping outputs valid.
    \item \textbf{Sparse ranking supervision}. A hard binary reward (1 for the correct item, 0 otherwise) offers little guidance on ranking quality. We introduce an auxiliary ranking-shaped reward to penalize harder negatives with lower scores. Additionally, dense signals such as semantic similarity and collaborative scores are explored to furnish richer supervision.
\end{itemize}
\subsubsection{Sampling Strategy}
\label{subsection:sampling}

A practical obstacle when porting RLVR to recommendation is the poor sampling diversity brought by the limited action space: querying the policy multiple times with the same prompt often returns identical items, so the model observes few distinct negatives. We measure diversity via
\[
\mathrm{Div}\!\left(\{e_k\}_{k=1}^G\right)=
\frac{\bigl|\text{Unique}\bigl(\{e_k\}_{1}^{G}\bigr)\bigr|}{G},
\]
where the numerator counts unique items among the $G$ generations. Higher values indicate richer supervision.

Two complementary remedies are investigated:

\begin{itemize}
    \item \textbf{Dynamic Sampling} \citep{DAPO}. We first over-sample, then pick a subset that (i) must include the ground-truth item and (ii) maximizes internal diversity. Although helpful, this demands extra forward passes and still deteriorates as training progresses.

    \item \textbf{Beam Search}. We ultimately switch to beam search without length normalization \citep{D3, ReRe}. By construction, all beams differ, so the method guarantees zero duplication within each group and yields better diversity–efficiency trade-offs.
\end{itemize}

Based on our findings, MiniOneRec ultimately employs constrained beam search as its default sampler, ensuring that every generated item is valid while still providing a diverse set of candidate trajectories.

\subsubsection{Reward Design}
\label{method:reward}

Recommendation models are usually judged with ranking measures such as NDCG. In contrast, the standard GRPO setup supplies a binary reward, giving a value of $1$ to the true item and $0$ to every other candidate. This strategy treats all negatives as equally harmful. Earlier studies \citep{SGL,sdpo} have shown that focusing on hard negatives produces stronger rankers. Motivated by these results, we introduce a rank-aware reward that assigns different penalties based on how prominently a negative item appears in the model’s own ranking \citep{ReRe}.

Given a negative candidate $e_k$ whose generation probability ranks $\rho_k$ (with $\rho=1$ the most probable), we set
\[
\tilde{R}_{\text{rank}}(e_k,e_t)=
\begin{cases}
0, & e_k=e_t,\\
-\dfrac{1}{\log(\rho_k+1)}, & \text{otherwise},
\end{cases}\;\;
R_{\text{rank}}(e_k,e_t)=
-\dfrac{\tilde{R}_{\text{rank}}(e_k,e_t)}
{\sum_{j=1}^{G}\tilde{R}_{\text{rank}}(e_j,e_t)}.
\]
Thus, negatives that the model was very confident about (low $\rho_k$) receive stronger penalties.

The final reward combines rule-based and ranking components:
\[
R(e_k,e_t)=R_{\text{rule}}(e_k,e_t)+R_{\text{rank}}(e_k,e_t),
\]
where the rule-based term is
\[
R_{\text{rule}}(e_k,e_t)=
\begin{cases}
1, & e_k = e_t,\\
0, & \text{otherwise}.
\end{cases}
\]
The mixed reward discussed combines binary correctness with a soft ranking term, helping the LLM tell hard negatives apart. Yet recommendation data contain additional, under-utilized cues. To see whether such information can enhance or even replace the rule-based component in RLVR, we experiment with the following collaborative reward. Specifically, for every item suggested by the policy, we obtain its logit from a pre-trained collaborative-filtering model and pass that score back as reward, thereby injecting knowledge extracted from historical user–item interactions. MiniOneRec ultimately adopts the combined ranking-and-rule reward as its default choice.




% \section{Methodology}
\label{method}

% \begin{figure}
%     \centering
%     \includegraphics[width=0.85\textwidth]{figures/MiniOneRec_framework.pdf}
%     \vspace{-10pt}
%     \caption{MiniOneRec framework. RQ-VAE builds the item SID codebook. We then perform SFT to warm up the LLM and obtain an initial alignment. In RL, beam search with constrained decoding, thereby the model sequentially produces a ranked list of distinct, valid SIDs. GRPO updates the policy, and SID alignment is enforced end-to-end. This alignment objective is preserved throughout both the SFT and RL stages, fostering deeper semantic understanding.}
%     \label{fig:framework}
%     \vspace{-10pt}
% \end{figure}


To address the key shortcomings of the existing SFT-then-RL pipeline for generative recommendation, we propose \textbf{MiniOneRec}.  
The framework elevates the performance ceiling by aligning the LLM with the SID space throughout the entire training process and applying preference optimization in the RL stage, as illustrated in Figure \ref{fig:framework}.

\subsection{Full-Process SID Alignment}

As demonstrated in Section~\ref{llmmatter}, aligning world knowledge with item SIDs is beneficial for generative recommendation.  
Hence, instead of the paradigm that trains only on SIDs \citep{TIGER, OneRec, mtgr}, our LLM-based recommender explicitly strengthens the link between language understanding and collaborative semantics by injecting a set of alignment objectives:
\begin{itemize}[leftmargin=*]
    \item \textbf{Semantic Tasks.} 
    Given a chronologically ordered sequence of historical SIDs and an explicit task instruction, the LLM is asked to predict the SID of the next item the user is likely to interact with. 
    \item \textbf{Alignment Tasks.} % Bridging tasks that transfer knowledge between the semantic and textual spaces. 
It comprises a series of alignment tasks between the textual space and the SID space. Through these tasks, we encourage a bidirectional mapping that grounds SIDs in language and injects linguistic knowledge into SID representations.
\end{itemize}
Representative tasks from each category are jointly optimized throughout the entire SFT-then-RL pipeline, enabling the LLM to fully exploit its world knowledge, deepen its understanding of SIDs, and ultimately boost generative-recommendation performance. During the RL phase, we employ constrained decoding, limiting the output space to a precompiled dictionary that contains the SID of each item as well as its canonical title. This restriction ensures that the LLM can emit only legal identifiers, making it straightforward to compute a rule-based, verifiable reward signal. Detailed examples of the prompts are provided in the Appendix \ref{appendix:align_prompt}.

% \subsection{SID-Natigated Reinforcement Learning}
\subsection{Preference Optimization Reinforcement Learning}
% To fully exploit the fine-grained signals carried by SIDs, we introduce SID-Natigated Reinforcement Learning (SIN-RL).  
% SIN-RL comprises two complementary components—curriculum sampling and reward modeling.  
% By leveraging the codebook’s inherent coarse-to-fine hierarchy, SIN-RL steers the agent toward harder examples in a structured manner, thereby improving its capability in the challenging regions of the data distribution.
\tjf{Adapting RLVR to recommendation is not straightforward, as it raises two key challenges in sampling and reward design.
First, unlike open-ended language generation, the output space of recommendation is constrained to valid item SIDs. 
It is much narrower and makes repeated sampling from the same prompt prone to severe duplication, which in turn leads to low sampling efficiency. 
To mitigate this, we incorporate dynamic sampling \citep{DAPO} and constrained beam search to improve sampling efficiency and expose the model to more informative negative items.

On the other hand, while ranking ability is essential for recommendation, the rule-based reward assigns zero to all negatives, providing only coarse-grained feedback. 
This exposes a broader challenge: reward design in recommendation remains under-explored, yet holds significant potential for richer supervision. 
To move beyond binary correctness, we propose an auxiliary ranking reward that penalizes harder negatives with lower scores, and further investigate the potential of dense rewards such as semantic and collaborative signals.

\subsubsection{Sampling Strategy}
\label{subsection:sampling}
A key challenge in adapting RLVR to recommendation lies in the low sampling efficiency caused by the narrow item space: repeatedly sampling from the same prompt often yields duplicate items, reducing the diversity of sampled negatives. 
To quantify this, we define generation diversity as
\begin{align}
    \label{eq:diversity}
    \text{div}(\{e_k\}_{k=1}^G)=\frac{\text{uni}(\{e_k\}_{k=1}^G)}{G},
\end{align}
where $\text{uni}(\{e_k\}_{k=1}^G)$ counts the number of unique items among $G$ generations. 
Higher diversity exposes the model to richer negative signals and thus improves ranking supervision.
To enhance sampling efficiency, we explore two complementary strategies.
We first explore dynamic sampling, which over-generates candidates and then selects a subset that balances inclusion of the target item with maximizing negative diversity. 
However, this strategy requires generating substantially more samples, and its diversity also inevitably declines as training progresses. 
Motivated by these limitations, we ultimately adopt beam search as our sampling method. 
Beam search ensures that all sampled items are distinct, thereby improving efficiency and diversity. 
Following \citep{Decodingmatters}, we remove length normalization to avoid amplification bias.
%, with additional details reported in Appendix~\ref{Appendix:ConsBeamSearch}.}


\subsubsection{Reward Design}
\label{method:reward}
Recommendation quality is usually measured by ranking metrics like NDCG, highlighting the importance of fine-grained ranking ability. 
However, the rule-based reward in GRPO only distinguishes the target item from all others, assigning $1$ to the positive and $0$ to every negative.
This implicitly assumes that all negatives contribute equally during optimization and only injects coarse-grained ranking information into LLMs.
% To further corroborate our analysis, we provide a gradient study of GRPO in Appendix \ref{Appendix:gradient}.

% assigning larger gradients to hard negatives 
Prior works have shown that emphasizing hard negatives improves the discriminative ability of recommendation models \citep{SGL, sdpo}.
Motivated by this insight, we introduce a ranking reward penalizing harder negatives more heavily. 
Specifically, for a generated negative item $e_k$, we compute its generation rank $\rho_k$ and assign a lower reward to negatives with higher-probability:
\begin{align}
    \hat{R}_\text{rank}(e_k, e_t)=\begin{cases} 
    0, & e_k=e_t \\
    -\frac{1}{\log(\rho_k+1)}, & \text{otherwise} 
    \end{cases},\\
    R_\text{rank}(e_k,e_t)=-\frac{\hat{R}_\text{rank}(e_k,e_t)}{\sum_{j=1}^G \hat{R}_\text{rank}(e_j,e_t)}. \label{eq:ranking_reward}
\end{align}
Finally, the overall reward combines the standard rule-based reward and the ranking reward:
\begin{align}
    R(e_k,e_t)=R_\text{rule}(e_k,e_t)+R_\text{rank}(e_k,e_t).
\end{align}
This combined reward integrates binary correctness with ranking signals, providing richer supervision that enhances the model’s ability to discriminate among negatives and improves ranking performance.
While the ranking reward provides finer-grained supervision than the binary rule-based reward, reward design in recommendation remains largely under-explored. 
Beyond binary correctness, recommendation scenarios naturally allow the use of richer dense signals that may further enhance alignment. 
\tjf{For further exploration}, we investigate two recommendation-specific dense rewards: 
\begin{itemize}[leftmargin=*]
    \item \textbf{Semantic reward.} This reward measures the semantic similarity between the generated item and the target item, encouraging LLMs to capture semantic relatedness between items. 
    \item \textbf{Collaborative reward.} To inject collaborative signals, each generated item is assigned the corresponding logit output by a traditional recommender system, reflecting collaborative filtering knowledge from historical user-item interactions.
\end{itemize}
These dense rewards offer a complementary perspective to binary correctness and allow us to examine whether richer supervision can 
substitute rule-based signals in RLVR for recommendation.}

% \subsubsection{SID-Prefix Curriculum Sampling}

% We first optimize the sampling strategy to preserve rollout diversity in RL for generative recommendation. In conventional LLM–RL training, LLM first applies dynamic sampling to obtain several candidate outputs and then computes a group–level reward.  
% Recent studies, however, report a rapid entropy drop during RL, which harms diversity \citep{entropy1, entropy2, entropy3, entropy4}.
% Inspired by these findings, we use a new sampling strategy tailored to generative recommendation so as to keep rollouts diverse.
% Specifically, when $k$ trajectories are required, we do not run $k$ independent dynamic samplings, each predicting a single item.  
% Instead, we execute one beam–search pass and take the top–$k$ item SIDs as the trajectory set for subsequent reward estimation.  
% Following prior work \citep{D3}, we remove length normalization in beam search to avoid bias amplification in the LLM-based recommender.

% Furthermore, we label a training instance $(x, y^{*})$ as difficult when none of the $k$ items generated in the current rollout matches the ground-truth SIDs $y^{*}$.  
% Once normal and hard tasks are separated, every hard sample $(x, y^*) \in \mathcal{D}_{\mathrm{sin}}$ is handled with a SID-Prefix curriculum schedule. At the beginning of training, we expose a long SID prefix; as learning progresses, the length of this prefix is gradually shortened, encouraging the model to explore on its own. Formally, we use p(t) to control the length of the SID-prefix:
%\begin{equation}
%     p(t) = 1- \frac{t}{T}
%\end{equation}
%where $T$ is the total number of RL steps and $t$ is the current step. Therefore, the value p(t) decreases linearly from 1 to 0, so the depth of the prefix guidance diminishes accordingly throughout training.

% For every hard instance $(x, y^{*}) \in \mathcal{D}_{\mathrm{sin}}$, we first compute the length of the SID-prefix, $L_{\mathrm{guide}}$, according to the schedule $p(t)$:
% \begin{equation}
%     L_{guide} = \lfloor p(t) \cdot L \rfloor 
% \end{equation}
% where
% $M$ denotes the number of hierarchical levels in the semantic codebook and
% $\lfloor \cdot \rfloor$ is the floor operator.
% We then truncate $y^{*}$ to its first $L_{\mathrm{guide}}$ tokens, denoted by $y_{L_{\mathrm{guide}}}^{*}$, and concatenate it with the original input $x$ to form a new prompt:
% \begin{equation}
%     x_{\mathrm{sin}} = x \oplus y_{L_{guide}}^*
% \end{equation}
% Finally, the policy produces a continuation
% $y \sim \pi_{\theta}(\,\cdot \mid x_{\mathrm{sin}})$,  
% upon which the RL reward is evaluated. % In summary, SID-Prefix curriculum sampling equips the policy network with adaptive SID guidance: a hierarchical prefix is injected at the early stage to alleviate dependence on the initial model’s capability and is gradually withdrawn as the model gains confidence. Through this curriculum, the LLM learns to handle hard cases autonomously, thereby pushing the upper bound of generative recommendation performance.

% \subsubsection{SID-level Reward Modeling}

% For a sampled item $e_i$ generated by the model, the naive rule-based reward $R_{\text{acc}}$ follows a binary scheme.
% The ground-truth item $e^{\text{pos}}$ is assigned 1, whereas every other candidate receives 0, as shown below: 
% \begin{equation}
% R_{\text{acc}}(e_i)=
% \begin{cases}
% 1, & e_i = e^{\text{pos}},\\[6pt]
% 0, & \text{otherwise}.
% \end{cases}
% \end{equation}
% Such sparsity treats all negative samples equally and thus fails to reflect their different levels of hardness.  
% We first follow the recent paradigm by constructing an auxiliary ranking score $R_{\text{rank}}$ that exploits the ordering information among candidate items.
% Specifically, a negative sample that appears higher in the generation list (i.e.\ is produced with a larger probability) should be penalized more. Let $p_i$ denote the position of a negative item $e_i^{\text{neg}}$. Its ranking reward is defined as the negative reciprocal of the natural logarithm of $(p_i+1)$, while the correct item $e^{\text{pos}}$ is given~0:
% \begin{equation}
% R_{\text{rank}}(e_i)=
% \begin{cases}
% 0, & e_i = e^{\text{pos}},\\[6pt]
% -\displaystyle\frac{1}{\log(p_i+1)}, & \text{otherwise}.
% \end{cases}
% \end{equation}
%
% Although the ranking reward supplies denser guidance, it still evaluates the whole sequence $\langle a \rangle \langle b \rangle \langle c \rangle $ as a single target and overlooks the multi-granular clues embedded in the SID codebook.  
% Consequently, for many difficult instances $(x, y^{*})\in \mathcal{D}_{\mathrm{sin}}$ the reward may still collapse to zero, leaving the sparsity issue unsolved.

% To address this limitation, we further view the generation of the hierarchical SID $\langle a \rangle \!\to\! \langle b \rangle \!\to\! \langle c \rangle$  as a step-by-step inference process and design SID-level reward.  
% Instead of relying on external Process-Reward Models, the proposed approach harnesses the codebook’s innate coarse-to-fine hierarchy to provide intermediate reasoning signals at each SID layer. This objective can be formally expressed as:
%
% \begin{equation}
% R_{\text{reason}}(e_i) = 1 - \beta^{k}
% \end{equation}
% where $k=f(e_i, e^{pos})$ denotes the deepest layer at which the generated sequence $e_i$ matches the ground-truth sequence $e^{pos}$, and $\beta \in (0,1)$ is a decay coefficient that modulates the reward increment across different layers and guarantees that the overall reward always lies in the interval $0 < R_{reason}(e_i) < 1$. The final total reward can be formally expressed as:
 
% \begin{equation}
% R(e_i)=
% \begin{cases}
% R_{\text{acc}}(e^i) + R_{\text{rank}}(e^i), & (x, y^*) \notin \mathcal{D}_{\mathrm{sin}},\\[6pt]
% R_{\text{reason}}(e^i), & \text{otherwise}.
% \end{cases}
% \end{equation}

% For challenging samples $(x, y^*) \in \mathcal{D}_{\mathrm{sin}}$, SID-Level Reward Modeling fully exploits SID’s hierarchical rewards to deliver fine-grained guidance.
% \section{Evaluation}

In this section, we first report the empirical performance of \textbf{MiniOneRec} on two real-world benchmarks \citep{amazon}, and compare it against a selected set of baselines covering conventional sequential recommenders, SID-based generative models, and recent LLM-powered recommenders.
In addition, we recast the SID-based next-item prediction task as
recommendation rule discovery and show that MiniOneRec offers extra gains
for cold-start users with scarce interactions and even for completely unseen domains.
We further conduct extensive ablation studies to pinpoint the components that most contribute to MiniOneRec’s effectiveness. 
In short, this section is organized to answer the following research 
questions:
\begin{itemize}[leftmargin=*]
    \item \textbf{RQ1:} How does MiniOneRec perform in comparison to other baseline methods?
    
    \item \textbf{RQ2:} How does MiniOneRec perform under completely unseen domains?
    
    \item \textbf{RQ3:} How do the designed components contribute to MiniOneRec's recommendation efficiency?
    % What is the impact of the designed components on the recommendation performance of MiniOneRec?
    \item \textbf{RQ4:} How does loss and recommendation metrics of MiniOneRec scale with different model size?
\end{itemize}

\textbf{Experimental Settings}
All conventional recommender baselines are optimized with binary cross-entropy (BCE) loss and the Adam optimizer.  
The learning rate is selected from $\{1\!\times\!10^{-2}, 1\!\times\!10^{-3}, 1\!\times\!10^{-4}\}$, while the weight-decay coefficient is tuned within $\{1\!\times\!10^{-2}, 1\!\times\!10^{-3}, 1\!\times\!10^{-4}, 1\!\times\!10^{-5}, 1\!\times\!10^{-6}\}$.  
A mini-batch size of 1024 is used throughout.  
For TIGER, we adopt T5 \citep{T5} as the encoder–decoder backbone and use Qwen3-Embedding-4B to generate item embedding.  
Every LLM-based method, including ours, is built upon Qwen2.5-Instruct-0.5B \citep{qwen2.5} to keep the computational footprint modest, and is trained with the AdamW optimizer.  
The SFT and preference-alignment data are processed in batches of 128, whereas RL batches contain 512 samples.  
We set the learning rate to $3\!\times\!10^{-4}$ for SFT and to $1\!\times\!10^{-5}$ for both S-DPO and MiniOneRec, together with a cosine decay scheduler.  
SFT runs for ten epochs with early stopping (patience = 1). S-DPO is trained for a single epoch, and we fix $\beta=0.1$ and sample three negative items. For D$^{3}$, the interpolation coefficient $\alpha$ is chosen from $\{0.8, 0.9, 1.0\}$.  
For the SID generation stage of MiniOneRec, we utilise Qwen3-Embedding-4B as the LLM encoder to transform item titles and descriptions into their corresponding embeddings. The tokenizer is trained on 8 GPUs with a per-device batch size of 2048. RQ-KMeans is trained layer-wise for 1\,000 steps per layer, with a learning rate of $1\!\times\!10^{-3}$. After embedding training, the SFT stage is conducted with a batch size of 128 for up to ten epochs (early stopping, patience = 1), followed by SID-Navigated RL for one epoch under the same $\beta$ and candidate settings as above.

\textbf{Datasets and Metrics.} 
We conduct extensive experiments on two real–world subsets of the Amazon Review corpus—\textit{Office} and \textit{Industrial}.  
% Among them, \textit{Office} and \textit{Industrial} are relatively small, each containing on average about 30\,000 user interaction sequences, whereas \textit{Books} is a large–scale dataset with more than 3\,million sequences.  
Following common practice, we adopt Hit Rate (HR@K) and Normalized Discounted Cumulative Gain (NDCG@K) to evaluate the top–$K$ recommendation accuracy. 

\textbf{Baselines.} 
Our baselines contain three categories: (1) Traditional recommendation models, including GRU4Rec \citep{GRU4Rec}, Caser \citep{Caser}, SASRec \citep{SASRec}; (2) Generative recommendation models: HSTU \citep{HSTU}, TIGER \citep{TIGER}, LC-Rec \citep{LCRec}; (3) LLM-based recommendation models, including BigRec \citep{BIGRec}, D$^3$ \citep{D3}, S-DPO \citep{sdpo}.


% Please add the following required packages to your document preamble:
% \usepackage{multirow}
\begin{table}[t]
\caption{Performance of MiniOneRec Compared to Traditional Methods, Generative Methods, and LLM-based Methods} 
\label{tab:overall}
\centering
\resizebox{0.9\textwidth}{!}{% 按比例缩放到页宽
\renewcommand{\arraystretch}{1.1}
\begin{tabular}{clllllll}
\hline
\multicolumn{1}{l}{Dateset}           & Methods    & HR@3   & NDCG@3 & HR@5   & NDCG@5 & HR@10  & NDCG@10 \\ \hline
\multirow{14}{*}{\textbf{Industrial}} & \multicolumn{7}{c}{\textbf{Traditional}}                          \\ \cline{2-8} 
                                      & GRU4Rec    & 0.0638 & 0.0542 & 0.0774 & 0.0598 & 0.0999 & 0.0669  \\
                                      & Caser      & 0.0618 & 0.0514 & 0.0717 & 0.0555 & 0.0942 & 0.0628  \\
                                      & SASRec     & 0.0790 & 0.0700 & 0.0909 & 0.0748 & 0.1088 & 0.0806  \\ \cline{2-8} 
                                      & \multicolumn{7}{c}{\textbf{Generative}}                           \\ \cline{2-8} 
                                      & HSTU       & 0.0927 & 0.0885 & 0.1037 & 0.0918 & 0.1163 & 0.0958  \\
                                      & TIGER      & 0.0852 & 0.0742 & 0.1010 & 0.0807 & 0.1321 & 0.0908  \\
                                      & LCRec      & 0.0915 & 0.0805 & 0.1057 & 0.0862 & 0.1332 & 0.0952  \\ \cline{2-8} 
                                      & \multicolumn{7}{c}{\textbf{LLM-based}}                            \\ \cline{2-8} 
                                      & BIGRec     & 0.0931 & 0.0841 & 0.1092 & 0.0907 & 0.1370 & 0.0997  \\
                                      & D$^3$         & 0.1024 & \uline{0.0991} & 0.1213 & 0.0989 & 0.1500 & \uline{0.1082}  \\
                                      & S-DPO      & \uline{0.1032} & 0.0906 & \uline{0.1238} & \uline{0.0991} & \uline{0.1524} & \uline{0.1082}  \\ \cline{2-8} 
                                      & \multicolumn{7}{c}{\textbf{Ours}}                                 \\ \cline{2-8} 
                                      & MiniOneRec      & \textbf{0.1125} & \textbf{0.0988} & \textbf{0.1255} & \textbf{0.1042} & \textbf{0.1508} & \textbf{0.1123}  \\
                                      \hline
\multirow{14}{*}{\textbf{Office}}     & \multicolumn{7}{c}{\textbf{Traditional}}                          \\ \cline{2-8} 
                                      & GRU4Rec    & 0.0629 & 0.0528 & 0.0789 & 0.0595 & 0.1019 & 0.0669  \\
                                      & Caser      & 0.0748 & 0.0615 & 0.0865 & 0.0664 & 0.1093 & 0.0737  \\
                                      & SASRec     & 0.0861 & 0.0769 & 0.0949 & 0.0805 & 0.1120 & 0.0858  \\ \cline{2-8} 
                                      & \multicolumn{7}{c}{\textbf{Generative}}                           \\ \cline{2-8} 
                                      & HSTU       & 0.1134 & 0.1031 & 0.1252 & 0.1079 & 0.1400 & 0.1126  \\
                                      & TIGER      & 0.0986 & 0.0852 & 0.1163 & 0.0960 & 0.1408 & 0.1002  \\
                                      & LCRec      & 0.0921 & 0.0807 &    0.1048 & 0.0859 & 0.1237 & 0.0920  \\ \cline{2-8} 
                                      & \multicolumn{7}{c}{\textbf{LLM-based}}                            \\ \cline{2-8} 
                                      & BIGRec     & 0.1069 & 0.0961 & 0.1204 & 0.1017 & 0.1434 & 0.1091  \\
                                      & D$^{3}$         & \uline{0.1204} & \uline{0.1055} & \uline{0.1406} & \uline{0.1139} & \uline{0.1634} & 0.1213  \\
                                      & S-DPO      & 0.1169 & 0.1033 & 0.1356 & 0.1110 & 0.1587 & \textbf{0.1255}  \\ \cline{2-8} 
                                      & \multicolumn{7}{c}{\textbf{Ours}}                                 \\ \cline{2-8} 
                                      & MiniOneRec      & \textbf{0.1217} & \textbf{0.1088} & \textbf{0.1420} & \textbf{0.1172} & \textbf{0.1634} & \uline{0.1242}  \\
                                      \hline
\end{tabular}
}
\end{table}

\subsection{Performance Comparison (RQ1)}

We conduct a comprehensive evaluation of \textsc{MiniOneRec} on three benchmark datasets—Industrial, Toys, and Books; the results are summarized in Table \ref{tab:overall}. Two major observations emerge:

\begin{itemize}[leftmargin=*]
    \item \textbf{Utility of LLM World Knowledge.} LLM-based recommenders such as BIGRec and D$^3$ markedly outperform classical paradigms like GRU4Rec and Caser, confirming that injecting the broad world knowledge encoded in LLMs can substantially boost recommendation quality.
    \item \textbf{Effectiveness of MiniOneRec.} By incorporating fine-grained SID information into the RL loop and aligning the entire generation trajectory with the task objective, MiniOneRec delivers a marked improvement over existing generative recommendation approaches, achieving consistent gains across all evaluated metrics. When contrasted with the current SOTA LLM-based recommendation paradigms, it continues to outperform the strongest baselines. Moreover, by operating in the compact SID space rather than on verbose textual titles, MiniOneRec requires substantially fewer context tokens and yields faster inference, translating into lower latency and smaller memory footprints at serving time.
\end{itemize}

% \begin{table}[t]
% \caption{Performance of MiniOneRec and its variants on completely unseen recommendation domains} 
% \label{tab:ood}
% \resizebox{\textwidth}{!}{% 按比例缩放到页宽
% \begin{tabular}{clllllll}
% \hline
% \multicolumn{1}{l}{Dataset}      & Method                    & HR@3            & NDCG@3          & HR@5            & NDCG@5          & HR@10           & NDCG@10         \\ \hline
% \multirow{5}{*}{\textbf{Office}} & GRU4Rec                   & 0.0629          & 0.0528          & 0.0789          & 0.0595          & 0.1019          & 0.0669          \\
%                                  & Qwen-Text               & 0.0031          & 0.0021          & 0.0044          & 0.0026          & 0.0057          & 0.0030          \\
%                                  & Qwen-SID                 & 0.0300          & 0.0214          & 0.0456          & 0.0282          & 0.0733          & 0.0373          \\
%                                  % & MiniOneRec & 0.0068          & 0.0056          & 0.0080          & 0.0061          & 0.0138          & 0.0080          \\
%                                  & MiniOneRec-w/o SFT     & \uline{0.0553} & \uline{0.0433} & \uline{0.0691} & \uline{0.0489} & \uline{0.0892} & \uline{0.0553} \\ \hline
% \end{tabular}
% }
% \end{table}

\subsection{Unseen-Domain Performance Evaluation (RQ2)}

To assess MiniOneRec's generalization to out-of-distribution (OOD) data, we conduct an unseen-domain study termed \emph{SID pattern discovery}. 
Specifically, the model is trained on the source domain Industrial and evaluated on a completely unseen target domain Office. Given that prior studies have shown SFT can overfit to the training domain and degrade OOD performance \citep{RLvsSFT, RLReallyLearn, SFT-then-RL, RevisitingRL}, we introduce an RL-only variant, MiniOneRec-w/ RL, specifically to provide a version focused on OOD generalization. We benchmark three systems:  
(1) GRU4Rec, trained and tested on Office.
(2) Qwen-Text, which encodes the user’s interaction history as plain text and predicts the next item in textual form;  
(3) Qwen-SID, which represents the same history as a sequence of SIDs and predicts the next SID;  
(4) MiniOneRec-w/ RL, which is trained solely with RL on Industrial and is tested on Office without SFT.  

As reported in Table~\ref{tab:ood}, Qwen-Text performs rather poorly, whereas Qwen-SID is markedly better, demonstrating that the structured SID space is easier for a language model to exploit. 
Although MiniOneRec-w/ RL lags behind the full MiniOneRec on in-domain metrics, its RL-only optimization offers strong transferability, yielding competitive accuracy on the unseen Office catalogue surprisingly.
Despite the substantial domain shift and the possible semantic mismatch among SIDs, MiniOneRec can still discover transferable interaction patterns and produce high-quality recommendations for a brand-new catalogue, highlighting the encouraging unseen-domain potential of our framework.



% \subsection{Training and Inference Cost Comparison (RQ3)}

% To comprehensively evaluate the resource efficiency of MiniOneRec, we compare its training cost with that of SOTA LLM-based recommenders fine-tuned via SFT. Specifically, we report GPU-hours to provide a more granular view of the overall expenditure.
% We carry out the cost analysis on the industry scale Books dataset, as illustrated in Figure x. Although MiniOneRec still adopts a two-stage pipeline of SFT followed by RL, its inputs and outputs are encoded as SID sequences, whose token count is an order of magnitude smaller than the free form text used by previous LLM-based recommenders. On the Books corpus we reserve only a small held out split for the RL phase, so the additional computation remains marginal. Overall, MiniOneRec delivers marked performance gains without inflating the training budget and simultaneously preserves the multi-modal SID interface.
% Beyond training cost, we measure online efficiency by reporting average latency per recommendation request. Owing to its compact SID representation, MiniOneRec requires substantially fewer decoding steps during inference than other LLM-based baselines. As displayed in Figure x, MiniOneRec cuts inference latency by 70\%, while using the same GPU batch size. These results indicate that MiniOneRec is not only cost-effective during training but also practical for real-time serving scenarios.

\subsection{Ablation Study (RQ3)}

To validate the effectiveness of each component in the MiniOneRec framework, we compare it with the following alternative approaches. 

\subsubsection{Aligning Strategy}

\begin{table}[t]
\caption{Performance of MiniOneRec and its variants on completely unseen recommendation domains} 
\label{tab:ood}
\resizebox{\textwidth}{!}{% 按比例缩放到页宽
\begin{tabular}{clllllll}
\hline
\multicolumn{1}{l}{Dataset}      & Method                    & HR@3            & NDCG@3          & HR@5            & NDCG@5          & HR@10           & NDCG@10         \\ \hline
\multirow{5}{*}{\textbf{Office}} & GRU4Rec                   & 0.0629          & 0.0528          & 0.0789          & 0.0595          & 0.1019          & 0.0669          \\
                                 & Qwen-Text               & 0.0031          & 0.0021          & 0.0044          & 0.0026          & 0.0057          & 0.0030          \\
                                 & Qwen-SID                 & 0.0300          & 0.0214          & 0.0456          & 0.0282          & 0.0733          & 0.0373          \\
                                 % & MiniOneRec & 0.0068          & 0.0056          & 0.0080          & 0.0061          & 0.0138          & 0.0080          \\
                                 & MiniOneRec-w/o SFT     & \uline{0.0553} & \uline{0.0433} & \uline{0.0691} & \uline{0.0489} & \uline{0.0892} & \uline{0.0553} \\ \hline
\end{tabular}
}
\end{table}

We benchmark the full model against three carefully designed variants:
(1) \textsc{MiniOneRec}--\textsc{w/o~Align}: A pure \textit{SID$\rightarrow$SID} paradigm: the input consists of SID-organized user histories, and the target is the SID of the next item. No cross-modal alignment is applied in either stage. 
(2) \textsc{MiniOneRec}--\textsc{w/~SFTAlign}: Alignment tasks are used only during the SFT stage, whereas the RL stage is trained on SID-only data.
(3) \textsc{MiniOneRec}--\textsc{w/~RLAlign}: The SFT stage relies on SID-only supervision, while alignment tasks are introduced solely in the RL stage.
% \begin{itemize}[leftmargin=*]
%     \item \textsc{MiniOneRec}--\textsc{w/o~Align.} A pure \textit{SID$\rightarrow$SID} paradigm: the input consists of SID-organized user histories, and the target is the SID of the next item. No cross-modal alignment is applied in either stage. 
%     \item \textsc{MiniOneRec}--\textsc{w/~SFTAlign.} Alignment tasks are used only during the SFT stage, whereas the RL stage is trained on SID-only data. 
%     \item \textsc{MiniOneRec}--\textsc{w/~RLAlign.} The SFT stage relies on SID-only supervision, while alignment tasks are introduced solely in the RL stage.
% \end{itemize}

As illustrated in Figure \ref{fig:ab1}, the full MiniOneRec with full-process SID alignment achieves the best results across all metrics. The \textsc{MiniOneRec}--\textsc{w/o~Align} variant performs the worst, underscoring the importance of grounding world knowledge when generating SIDs. Notably, although introducing the alignment objective directly in the RL stage is highly challenging for an LLM that has not been pre-conditioned by SFT, \textsc{MiniOneRec}--\textsc{w/~RLAlign} still yields a non-trivial gain. We attribute this improvement to our SID-Guided RL scheme, which offers the agent additional opportunities to obtain valid rewards on hard examples that conventional RL would likely miss.

\subsubsection{Sampling Strategy}

We contrast the full model with three variants that differ only in the way trajectories are collected:
(1) \textsc{MiniOneRec}--\textsc{Common} that uses a conventional top-$k$ decoding scheme to generate the required number of trajectories.
(2) \textsc{MiniOneRec}--\textsc{Dynamic}, which implements our dynamic sampler that first produces $\tfrac32$ times the target trajectory budget and then keeps as many distinct items as possible for RL optimization.
(3) \textsc{MiniOneRec}--\textsc{w/~Beam} that retains only beam search; no SID-prefix guidance is applied, and all examples maintain their original difficulty level.
% \begin{itemize}[leftmargin=*]
%     \item \textsc{MiniOneRec}--\textsc{Common.} 
%     Uses a conventional top-$k$ decoding scheme to generate the required number of trajectories.
%     \item \textsc{MiniOneRec}--\textsc{Dynamic.} 
%     Implements our dynamic sampler that first produces $\tfrac32$ times the target trajectory budget and then keeps as many distinct items as possible for RL optimization.
%     \item \textsc{MiniOneRec}--\textsc{w/~Beam.} 
%     Retains only beam search; no SID-prefix guidance is applied, and all examples maintain their original difficulty level.
% \end{itemize}

As Figure \ref{fig:ab2} illustrates, the full MiniOneRec achieves the best overall performance. Moreover, the \textsc{MiniOneRec}--\textsc{w/~Beam} variant attains higher accuracy than \textsc{MiniOneRec}--\textsc{Dynamic} while requiring only two-thirds of its sampling budget, showing that beam search is a more cost-effective backbone. These findings motivate our final design that merges beam search with SID-prefix guidance.

\subsubsection{Reward Design}

Three variants are compared:
(1) \textsc{MiniOneRec}--\textsc{w/~acc} that uses the accuracy reward only;
(2) \textsc{MiniOneRec}--\textsc{w/o~SID} that removes the SID-level reward for hard cases; the LLM is optimized with the accuracy plus rank rewards.
(3) \textsc{MiniOneRec}--\textsc{w/~SIDAll} that applies the SID-level reward to every training sample rather than restricting it to difficult ones.

% \begin{itemize}[leftmargin=*]
%    \item \textsc{MiniOneRec}--\textsc{w/~acc.} 
%    Uses the accuracy reward only.
%    \item \textsc{MiniOneRec}--\textsc{w/o~SID.} Removes the SID-level reward for hard cases; the LLM is optimized with the accuracy plus rank rewards.
%    \item \textsc{MiniOneRec}--\textsc{w/~SIDAll.} 
%    Applies the SID-level reward to every training sample rather than restricting it to difficult ones.
% \end{itemize}
As shown in the Figure \ref{fig:ab3}, the full model, which deploys SID-level reward only for hard samples, achieves the best overall performance. To be noticed, rewarding all samples at the SID-level slightly degrades performance. We hypothesize that, for easy instances already covered by the accuracy and ranking rewards, an additional SID-level signal may dilute the optimization focus and introduce a mismatch with the evaluation metrics (HR@K and NDCG@K). For hard cases, however, the hierarchical structure encoded in SIDs provides intermediate guidance where conventional rewards are sparse, steering the model toward correct reasoning paths.

\begin{figure*}[t]
  \centering
  \begin{subfigure}[t]{0.28\textwidth}
    \centering
    \includegraphics[width=\textwidth]{figures/ab1.png}
    \caption{Aligning Strategy.}
    \label{fig:ab1}
\end{subfigure}\hspace{2mm}
\begin{subfigure}[t]{0.28\textwidth}
    \centering
        \includegraphics[width=\textwidth]{figures/ab2.png}
    \caption{Sampling Strategy.}
    \label{fig:ab2}
\end{subfigure}
\centering
  \begin{subfigure}[t]{0.28\textwidth}
    \centering
    \includegraphics[width=\textwidth]{figures/ab3.png}
    \caption{Reward Design.}
    \label{fig:ab3}
\end{subfigure}\hspace{2mm}
  \caption{Study on the effectiveness of MiniOneRec’s individual components.
Figure \ref{fig:ab1} examines model performance under different alignment strategies;
Figure \ref{fig:ab2} investigates various sampling strategies;
Figure \ref{fig:ab3} evaluates the impact of alternative reward designs.}
  \label{fig:ablation}
    \vspace{-10pt}
\end{figure*}

\subsection{Scaling Study (RQ4)}
\section{Training}
\label{training_detail}

MiniOneRec first converts item text into discrete codes. Titles and descriptions are embedded by the Qwen3-Embedding-4B encoder, after which RQ-VAE performs residual quantization. This tokenizer is trained on a single GPU with a batch size of 20\,480, a learning rate of $1\times10^{-3}$, and 10,000 training epochs. The resulting SID vocabulary is plugged into a Qwen2.5-Instruct backbone. SFT takes place on eight NVIDIA H100, each holding 128 samples. Training lasts up to ten epochs with early stopping (patience one epoch). The initial learning rate is $3\times10^{-4}$ and follows a cosine decay schedule. Starting from the SFT checkpoint, we apply GRPO for two additional epochs while keeping the KL weight $\beta$ unchanged. During roll-out, beam search with width 16 is adopted, so every input generates sixteen distinct candidate sequences.

In the performance comparison with existing recommenders, conventional recommenders are trained with binary cross-entropy and the Adam optimizer; learning rates are chosen from \{1e-2, 1e-3, 1e-4\}, and the weight-decay term is scanned over \{1e-2, 1e-3, 1e-4, 1e-5, 1e-6\}. The batch size is 1\,024.  
For TIGER, we rely on a T5 encoder–decoder and reuse Qwen3-Embedding-4B for item embeddings.  
All LLM-powered systems, including ours, share the Qwen2.5-Instruct backbone and use AdamW. Batches contain 128 examples for SFT and preference alignment, and 512 samples for RL. The learning rate is $3\times10^{-4}$ during SFT, while S-DPO and RL use $1\times10^{-5}$. S-DPO runs for one epoch with $\beta=0.1$ and samples three negative items; D$^{3}$ tries interpolation factors $\alpha$ of 0.8, 0.9, and 1.0.

\section{Evaluation}
\label{results}

This section details our empirical study. We begin with a scaling investigation on two real-world benchmarks \citep{amazon}, examining how MiniOneRec’s loss curves change as model size grows. We then benchmark MiniOneRec against three groups of baselines: traditional sequential recommenders, SID-based generators, and recent LLM-driven systems.  
To evaluate the model’s capability for cross-domain recommendation generalization, we frame SID next-item prediction as a task of recommendation rule discovery and evaluate the model on domains that never appeared during training.  
Comprehensive ablation tests follow, which help us locate the parts of the framework that contribute most to the final score. Finally, we look into how the broad knowledge stored in pre-trained LLM weights affects generative recommendation.
In summary, we study the following questions within this section:

\begin{itemize}
    \item How does the loss of MiniOneRec scale with different model size?
    
    \item How does MiniOneRec perform in comparison to other baseline methods?
    
    \item How does MiniOneRec perform under completely unseen domains?
    
    \item How do the designed components contribute to MiniOneRec's recommendation efficiency?

    \item How do the pre-trained weights of the LLM influence the performance of generative recommendation?
    % What is the impact of the designed components on the recommendation performance of MiniOneRec?

\end{itemize}

\subsection{Evaluation Setup}

% \textbf{Datasets and Metrics.} 

Experiments are carried out on two real-world slices of the Amazon Review dataset \citep{amazon}, namely \textit{Office} and \textit{Industrial}. To measure top-$K$ recommendation accuracy, we follow standard practice and compute Hit Rate (HR@K) as well as Normalized Discounted Cumulative Gain (NDCG@K).


% Please add the following required packages to your document preamble:
% \usepackage{multirow}
\begin{table}[t]
\caption{Performance of MiniOneRec Compared to Traditional Methods, Generative Methods, and LLM-based Methods} 
\label{tab:overall}
\centering
\begin{adjustbox}{max width=0.9\textwidth}
\renewcommand{\arraystretch}{1.00}
% \begin{tabular}{clllllll}
\begin{tabular}{cccccccc}
\toprule
\multicolumn{1}{c}{Datasets}            & Methods    & HR@3   & NDCG@3 & HR@5   & NDCG@5 & HR@10  & NDCG@10 \\ \midrule
% \multicolumn{1}{l}{Datesets}           & Methods    & HR@3   & NDCG@3 & HR@5   & NDCG@5 & HR@10  & NDCG@10 \\ \midrule
\multirow{14}{*}{\textbf{Industrial}} & \multicolumn{7}{c}{\textbf{Traditional}}                          \\ \cmidrule(lr){2-8} 
                                      & GRU4Rec    & 0.0638 & 0.0542 & 0.0774 & 0.0598 & 0.0999 & 0.0669  \\
                                      & Caser      & 0.0618 & 0.0514 & 0.0717 & 0.0555 & 0.0942 & 0.0628  \\
                                      & SASRec     & 0.0790 & 0.0700 & 0.0909 & 0.0748 & 0.1088 & 0.0806  \\ \cmidrule(lr){2-8}
                                      & \multicolumn{7}{c}{\textbf{Generative}}                           \\ \cmidrule(lr){2-8}
                                      & HSTU       & 0.0927 & 0.0885 & 0.1037 & 0.0918 & 0.1163 & 0.0958  \\
                                      & TIGER      & 0.0852 & 0.0742 & 0.1010 & 0.0807 & 0.1321 & 0.0908  \\
                                      & LCRec      & 0.0915 & 0.0805 & 0.1057 & 0.0862 & 0.1332 & 0.0952  \\ \cmidrule(lr){2-8}
                                      & \multicolumn{7}{c}{\textbf{LLM-based}}                            \\ \cmidrule(lr){2-8}
                                      & BIGRec     & 0.0931 & 0.0841 & 0.1092 & 0.0907 & 0.1370 & 0.0997  \\
                                      & D$^3$         & 0.1024 & \uline{0.0991} & 0.1213 & 0.0989 & 0.1500 & \uline{0.1082} \\
                                      & S-DPO      & \uline{0.1032} & 0.0906 & \uline{0.1238} & \uline{0.0991} & \uline{0.1524} & \uline{0.1082}  \\ \cmidrule(lr){2-8}
                                      & \multicolumn{7}{c}{\textbf{Ours}}                                 \\ \cmidrule(lr){2-8}
                                      & \cellcolor[HTML]{ECF4FF}\textbf{MiniOneRec}      & \cellcolor[HTML]{ECF4FF}\textbf{0.1143} & \cellcolor[HTML]{ECF4FF}\textbf{0.1011} & \cellcolor[HTML]{ECF4FF}\textbf{0.1321} & \cellcolor[HTML]{ECF4FF}\textbf{0.1084} & \cellcolor[HTML]{ECF4FF}\textbf{0.1586} & \cellcolor[HTML]{ECF4FF}\textbf{0.1167}  \\
                                      \midrule
\multirow{14}{*}{\textbf{Office}}     & \multicolumn{7}{c}{\textbf{Traditional}}                          \\ \cmidrule(lr){2-8}
                                      & GRU4Rec    & 0.0629 & 0.0528 & 0.0789 & 0.0595 & 0.1019 & 0.0669  \\
                                      & Caser      & 0.0748 & 0.0615 & 0.0865 & 0.0664 & 0.1093 & 0.0737  \\
                                      & SASRec     & 0.0861 & 0.0769 & 0.0949 & 0.0805 & 0.1120 & 0.0858  \\ \cmidrule(lr){2-8}
                                      & \multicolumn{7}{c}{\textbf{Generative}}                           \\ \cmidrule(lr){2-8}
                                      & HSTU       & 0.1134 & 0.1031 & 0.1252 & 0.1079 & 0.1400 & 0.1126  \\
                                      & TIGER      & 0.0986 & 0.0852 & 0.1163 & 0.0960 & 0.1408 & 0.1002  \\
                                      & LCRec      & 0.0921 & 0.0807 &    0.1048 & 0.0859 & 0.1237 & 0.0920  \\ \cmidrule(lr){2-8}
                                      & \multicolumn{7}{c}{\textbf{LLM-based}}                            \\ \cmidrule(lr){2-8}
                                      & BIGRec     & 0.1069 & 0.0961 & 0.1204 & 0.1017 & 0.1434 & 0.1091  \\
                                      & D$^{3}$         & \uline{0.1204} & \uline{0.1055} & \uline{0.1406} & \uline{0.1139} & \textbf{0.1634} & 0.1213  \\
                                      & S-DPO      & 0.1169 & 0.1033 & 0.1356 & 0.1110 & \uline{0.1587} & \textbf{0.1255}  \\ \cmidrule(lr){2-8}
                                      & \multicolumn{7}{c}{\textbf{Ours}}                                 \\ \cmidrule(lr){2-8}
                                      & \cellcolor[HTML]{ECF4FF}\textbf{MiniOneRec}      & \cellcolor[HTML]{ECF4FF}\textbf{0.1217} & \cellcolor[HTML]{ECF4FF}\textbf{0.1088} & \cellcolor[HTML]{ECF4FF}\textbf{0.1420} & \cellcolor[HTML]{ECF4FF}\textbf{0.1172} & \cellcolor[HTML]{ECF4FF}\textbf{0.1634} & \cellcolor[HTML]{ECF4FF}\uline{0.1242}  \\
                                      \bottomrule
\end{tabular}
\end{adjustbox}
\end{table}

\subsection{Scaling}
We demonstrate the scaling capabilities of MiniOneRec, where the convergence loss decrease consistently as the model size increases. 
As Figure \ref{fig:scaling} shows, under the generative recommendation paradigm, there is a clear correlation between model size and loss. As the LLM scale increases, the convergence loss continues to decrease. Furthermore, Figure \ref{fig:scaling_eval_loss} tracks the evaluation loss on the SFT training set, recorded every 0.5 epoch. It is evident that models with larger parameter counts maintain lower evaluation losses throughout nearly the entire training process and converge more rapidly.
Such a superior scaling effect demonstrates the potential of generative recommenders as the next-generation recommendation models.

% \subsection{Performance Comparison (RQ1)}
\subsection{Performance Comparison}

% \textbf{Experimental Settings}.
% All conventional recommender baselines are optimized with binary cross-entropy (BCE) loss and the Adam optimizer. The learning rate is selected from $\{1\!\times\!10^{-2}, 1\!\times\!10^{-3}, 1\!\times\!10^{-4}\}$, while the weight-decay coefficient is tuned within $\{1\!\times\!10^{-2}, 1\!\times\!10^{-3}, 1\!\times\!10^{-4}, 1\!\times\!10^{-5}, 1\!\times\!10^{-6}\}$. A mini-batch size of 1024 is used throughout. For TIGER, we adopt T5 \citep{T5} as the encoder–decoder backbone and use Qwen3-Embedding-4B to generate item embedding. Every LLM-based method, including ours, is built upon Qwen2.5-Instruct-0.5B \citep{qwen2.5} to keep the computational footprint modest, and is trained with the AdamW optimizer. The SFT and preference-alignment data are processed in batches of 128, whereas RL batches contain 512 samples. We set the learning rate to $3\!\times\!10^{-4}$ for SFT and to $1\!\times\!10^{-5}$ for both S-DPO and MiniOneRec, together with a cosine decay scheduler. SFT runs for ten epochs with early stopping (patience = 1). S-DPO is trained for a single epoch, and we fix $\beta=0.1$ and sample three negative items. For D$^{3}$, the interpolation coefficient $\alpha$ is chosen from $\{0.8, 0.9, 1.0\}$.  

% \textbf{Baselines.} 

Our baselines contain three categories: (1) Traditional recommendation models, including GRU4Rec \citep{GRU4Rec}, Caser \citep{Caser}, SASRec \citep{SASRec}; (2) Generative recommendation models: HSTU \citep{HSTU}, TIGER \citep{TIGER}, LC-Rec \citep{LCRec}; (3) LLM-based recommendation models, including BigRec \citep{BIGRec}, D$^3$ \citep{D3}, S-DPO \citep{sdpo}.

We evaluate \textsc{MiniOneRec} on the \textit{Industrial} and \textit{Office} benchmarks and summarize the outcomes in Table~\ref{tab:overall}. Two key insights stand out:


\begin{itemize}
    \item \textbf{Utility of LLM World Knowledge.} Recommenders powered by LLMs, such as BIGRec and D$^{3}$, clearly surpass traditional systems like GRU4Rec and Caser, showing that the broad knowledge embedded in LLMs translates into better recommendation accuracy.  
    \item \textbf{Effectiveness of MiniOneRec.} MiniOneRec goes a step further: by aligning the full generation process with the task objective and using reinforced preference optimization during RL, it consistently outperforms previous generative solutions across most reported metrics. Moreover, by operating in the compact SID space rather than on verbose textual titles, MiniOneRec requires substantially fewer context tokens and yields faster inference, translating into lower latency and smaller memory footprints at serving time.
\end{itemize}

% \subsection{Unseen-Domain Performance Evaluation (RQ2)}
\subsection{Transferability}
\label{subsec: Transferability}

We evaluate the out-of-distribution (OOD) robustness of MiniOneRec through an experiment we call \textit{SID pattern discovery}. The model is trained exclusively on the \textit{Industrial} domain and then deployed, without any further tuning, to the never-seen \textit{Office} domain. Prior work \citep{RLvsSFT, RLReallyLearn, SFT-then-RL, RevisitingRL} suggests that SFT may overfit to the source domain and hurt transfer, so we add an RL-only variant, MiniOneRec-w/ RL-OOD, that skips SFT entirely to emphasize generalization.
\begin{wrapfigure}[14]{r}{0.48\columnwidth}   % 14 可根据需要微调
    \centering
    \vspace{2pt}                             % 细调上边距
    \includegraphics[width=\linewidth]{figures/scaling_eval_loss.png}
    \vspace{-16pt}                             % 细调下边距
    \caption{Evaluation loss vs.\ SFT training epoch}
    \label{fig:scaling_eval_loss}
\end{wrapfigure}
The comparison involves four systems:  
(1) GRU4Rec, trained and tested inside the \textit{Office} domain;  
(2) Qwen-Text, which represents user histories as plain sentences and predicts the next item by its title with no additional fine-tuning;  
(3) Qwen-SID, which encodes the same history with SID tokens and produces the next SID with no additional fine-tuning;  
(4) MiniOneRec-w/ RL-OOD, trained via GRPO on \textit{Industrial} only and evaluated on \textit{Office}, to evaluate its OOD (out-of-distribution) performance.

Table~\ref{tab:ood} reports the outcomes. Qwen-Text performs poorly, while Qwen-SID does noticeably better, indicating that a structured SID vocabulary is easier for an LLM to exploit. Although MiniOneRec-w/ RL-OOD falls short of the full MiniOneRec on in-domain scores, its reinforcement-only training grants excellent transfer, achieving competitive accuracy on the unseen catalog. Despite the substantial domain shift and possible semantic drift among SIDs, MiniOneRec successfully uncovers reusable interaction patterns, underscoring the framework’s promise for cross-domain recommendation.


% \subsection{Training and Inference Cost Comparison (RQ3)}

% To comprehensively evaluate the resource efficiency of MiniOneRec, we compare its training cost with that of SOTA LLM-based recommenders fine-tuned via SFT. Specifically, we report GPU-hours to provide a more granular view of the overall expenditure.
% We carry out the cost analysis on the industry scale Books dataset, as illustrated in Figure x. Although MiniOneRec still adopts a two-stage pipeline of SFT followed by RL, its inputs and outputs are encoded as SID sequences, whose token count is an order of magnitude smaller than the free form text used by previous LLM-based recommenders. On the Books corpus we reserve only a small held out split for the RL phase, so the additional computation remains marginal. Overall, MiniOneRec delivers marked performance gains without inflating the training budget and simultaneously preserves the multi-modal SID interface.
% Beyond training cost, we measure online efficiency by reporting average latency per recommendation request. Owing to its compact SID representation, MiniOneRec requires substantially fewer decoding steps during inference than other LLM-based baselines. As displayed in Figure x, MiniOneRec cuts inference latency by 70\%, while using the same GPU batch size. These results indicate that MiniOneRec is not only cost-effective during training but also practical for real-time serving scenarios.

% \subsection{Ablation Study (RQ3)}
\subsection{Ablation Study}

To validate the effectiveness of each component in the MiniOneRec framework, we compare it with the following alternative approaches. 

\subsubsection{Aligning Strategy}

% Please add the following required packages to your document preamble:
% \usepackage{multirow}
% \usepackage[table,xcdraw]{xcolor}
% Beamer presentation requires \usepackage{colortbl} instead of \usepackage[table,xcdraw]{xcolor}
\begin{table}[]
\caption{Performance of MiniOneRec and its variants on completely unseen recommendation domains} 
\label{tab:ood}
\resizebox{\textwidth}{!}{% 按比例缩放到页宽
% \begin{tabular}{clllllll}
\begin{tabular}{cccccccc}
\toprule
Dataset                           & Method                          & HR@3                           & NDCG@3                         & HR@5                           & NDCG@5                         & HR@10                          & NDCG@10                        \\ \hline
                                  & \cellcolor[HTML]{EFEFEF}GRU4Rec & \cellcolor[HTML]{EFEFEF}0.0629 & \cellcolor[HTML]{EFEFEF}0.0528 & \cellcolor[HTML]{EFEFEF}0.0789 & \cellcolor[HTML]{EFEFEF}0.0595 & \cellcolor[HTML]{EFEFEF}0.1019 & \cellcolor[HTML]{EFEFEF}0.0669 \\
                                  & Qwen-Text                       & 0.0031                         & 0.0021                         & 0.0044                         & 0.0026                         & 0.0057                         & 0.0030                         \\
                                  & Qwen-SID                        & 0.0300                         & 0.0214                         & 0.0456                         & 0.0282                         & 0.0733                         & 0.0373                         \\
\multirow{-4}{*}{\textbf{Office}} & \cellcolor[HTML]{ECF4FF}MiniOneRec-w/ RL-OOD              & \cellcolor[HTML]{ECF4FF}0.0553                         & \cellcolor[HTML]{ECF4FF}0.0433                         & \cellcolor[HTML]{ECF4FF}0.0691                         & \cellcolor[HTML]{ECF4FF}0.0489                         & \cellcolor[HTML]{ECF4FF}0.0892                         & \cellcolor[HTML]{ECF4FF}0.0553                         \\ \hline
\end{tabular}
}
\end{table}

To isolate the impact of each component, we compare the full MiniOneRec with three pared-down variants.  
(1) \textsc{MiniOneRec}--\textsc{w/o~Align}: removes any language–SID alignment and treats recommendation purely as a SID-to-SID task.  
(2) \textsc{MiniOneRec}--\textsc{w/~SFTAlign}: keeps the alignment objective during SFT stage only, while RL uses SID data alone.  
(3) \textsc{MiniOneRec}--\textsc{w/~RLAlign}: SFT relies solely on SID supervision, and the alignment tasks are introduced later in the RL stage.

Figure~\ref{fig:ab1} summarizes the findings. The complete MiniOneRec, which maintains alignment throughout the whole pipeline, delivers the highest scores on every metric. The \textsc{MiniOneRec}--\textsc{w/o~Align} performs worst, indicating that grounding SID generation in world knowledge is essential.

\subsubsection{Sampling Strategy}

We study how different roll-out methods affect MiniOneRec by switching only the trajectory generator while keeping everything else fixed:  
(1) \textsc{MiniOneRec}--\textsc{Common} relies on a plain Top-$k$ decoder to produce exactly the required number of paths.  
(2) \textsc{MiniOneRec}--\textsc{Dynamic} follows our two-step sampler: it first draws one and a half times the budget, then retains as many unique items as possible for RL.  
The full model adopts beam search with width 16.

As plotted in Figure~\ref{fig:ab2}, the complete MiniOneRec delivers the highest accuracy while using roughly two-thirds of the samples needed by the dynamic variant, demonstrating that beam search is the most cost-efficient choice among the tested strategies.

\subsubsection{Reward Design}

Three variants are compared:
(1) \textsc{MiniOneRec}--\textsc{w/~ACC} that relies solely on a binary correctness signal;  
(2) \textsc{MiniOneRec}--\textsc{w/~Collaborative} that replaces the ranking term with logits taken from a frozen SASRec model so as to supply collaborative cues;  

As shown in the Figure \ref{fig:ab3}, the full MiniOneRec achieves the best overall performance. To be noticed, injecting the collaborative reward information into the RL process instead led to significant degradation. We hypothesize that, this stems from reward hacking: as recommendation accuracy declines, the reward continues to increase, revealing a misalignment between this collaborative reward signal and the true objective.

\begin{figure*}[t]
  \centering
  \begin{subfigure}[t]{0.28\textwidth}
    \centering
    \includegraphics[width=\textwidth]{figures/ab1.png}
    \caption{Aligning Strategy.}
    \label{fig:ab1}
\end{subfigure}\hspace{2mm}
\begin{subfigure}[t]{0.28\textwidth}
    \centering
        \includegraphics[width=\textwidth]{figures/ab2.png}
    \caption{Sampling Strategy.}
    \label{fig:ab2}
\end{subfigure}
\centering
  \begin{subfigure}[t]{0.28\textwidth}
    \centering
    \includegraphics[width=\textwidth]{figures/ab3.png}
    \caption{Reward Design.}
    \label{fig:ab3}
\end{subfigure}\hspace{2mm}
  \caption{Study on the effectiveness of MiniOneRec’s individual components.
Figure \ref{fig:ab1} examines model performance under different alignment strategies;
Figure \ref{fig:ab2} investigates various sampling strategies;
Figure \ref{fig:ab3} evaluates the impact of alternative reward designs.}
  \label{fig:ablation}
    \vspace{-10pt}
\end{figure*}



% Please add the following required packages to your document preamble:
% \usepackage{multirow}
\begin{table}[]
\caption{Performance of MiniOneRec initialized with pre-trained weights versus random initialization across two domains} 
\label{tab:pretrain-impact}
\resizebox{\textwidth}{!}{% 按比例缩放到页宽
\begin{tabular}{clllllll}
\toprule
\multicolumn{1}{l}{Datesets}          & Methods            & HR@3            & NDCG@3          & HR@5            & NDCG@5          & HR@10           & NDCG@10         \\ \midrule
\multirow{2}{*}{\textbf{Industrial}} & MiniOneRec-scratch & 0.0757          & 0.0672          & 0.0891          & 0.0726          & 0.1134          & 0.0804          \\
                                     & \cellcolor[HTML]{ECF4FF}\textbf{MiniOneRec}         & \cellcolor[HTML]{ECF4FF}\textbf{0.1125} & \cellcolor[HTML]{ECF4FF}\textbf{0.0988} & \cellcolor[HTML]{ECF4FF}\textbf{0.1259} & \cellcolor[HTML]{ECF4FF}\textbf{0.1046} & \cellcolor[HTML]{ECF4FF}\textbf{0.1546} & \cellcolor[HTML]{ECF4FF}\textbf{0.1139} \\ \midrule
\multirow{2}{*}{\textbf{Office}}     & MiniOneRec-scratch & 0.0959          & 0.0855          & 0.1057          & 0.0896          & 0.1196          & 0.0941          \\
                                     & \cellcolor[HTML]{ECF4FF}\textbf{MiniOneRec}         & \cellcolor[HTML]{ECF4FF}\textbf{0.1217} & \cellcolor[HTML]{ECF4FF}\textbf{0.1088} & \cellcolor[HTML]{ECF4FF}\textbf{0.1420} & \cellcolor[HTML]{ECF4FF}\textbf{0.1172} & \cellcolor[HTML]{ECF4FF}\textbf{0.1634} & \cellcolor[HTML]{ECF4FF}\textbf{0.1242} \\ \bottomrule
\end{tabular}
}
\end{table}


\subsection{Pre-trained LLM Impact}
\label{llmmatter}
To isolate the impact of pre-trained LLMs on generative recommendation, we instantiate two variants of MiniOneRec: one initialized from a general-purpose pre-trained LLM and the other trained from scratch with random weights. As shown in Table \ref{tab:pretrain-impact}, Experiments on two benchmark datasets reveal a consistent pattern: the model that starts from pre-trained weights significantly outperforms its randomly initialized counterpart. We conjecture that (i) the general reasoning ability acquired during large-scale language pre-training allows the model to cast the next-SID prediction task as a problem of pattern discovery, as discussed in Section~\ref{subsec: Transferability}, and (ii) the factual knowledge already encoded in the LLM offers a head start in understanding the real-world semantics behind each SID, part of which can be transferred to the recommendation domain.
% \section{Related Work}
\subsection{Generative Recommendation}

In recent years, generative recommendation has attracted considerable attention in both academia and industry. This emerging paradigm, usually built on the Transformer architecture, reformulates recommendation as an end-to-end next-item generation task and thereby raises the performance ceiling of recommender systems.
Early work TIGER \citep{TIGER} employs residual quantization (RQ-VAE) \citep{rqvae} to convert text embeddings—extracted from an item’s title and description—into discrete semantic IDs, which are then used for next-item prediction. LC-Rec \citep{LCRec} aligns an LLM with semantic IDs via multi-task learning, enabling the model to “understand” the IDs and perform generative recommendation. Other studies investigate how to build better semantic IDs to enhance generation quality: RecForest \citep{recforest} applies hierarchical k-means clustering and treats the cluster indices as tokens, while EAGER \citep{eagar} and TokenRec \citep{TOKENRec} integrate semantic and collaborative signals directly into the tokenizer.

Very recently, generative recommendation has been rolled out at an industrial scale to address the drawbacks of traditional cascade systems. MTGR \citep{mtgr} keeps the original deep learning recommendation model (DLRM) features, introduces user-level compression, and speeds up both training and inference for large-scale deployment. OneRec \citep{OneRec} lowers the serving cost with a Lazy Decoder-Only Architecture and stabilises training through an improved reinforcement learning algorithm.

\subsection{LLM and RL}
Reinforcement Learning (RL) \citep{RLsurvey, RLsurvey2} operates through the interaction mechanism between agents and the environment, rewarding correct actions with the core objective of maximizing cumulative returns. In recent years, researchers have introduced RL into the fine-tuning process of LLMs. Among these approaches, Reinforcement Learning with Human Feedback (RLHF) \citep{RLHF} typically employs the Proximal Policy Optimization (PPO) \citep{PPO} algorithm to achieve alignment with human preferences. However, the PPO algorithm faces technical challenges related to excessive memory consumption. To simplify the application of RL in LLMs, researchers have proposed innovative methods such as Direct Preference Optimization (DPO) \citep{DPO}. In the context of recommender systems, the s-DPO method \citep{sdpo} combines DPO with the softmax loss function, enhancing LLMs' capability in mining hard negative samples for recommender systems. Although these methods have improved efficiency, they still face challenges such as off-policy issues and performance ceilings that fall short of online methods. The recently proposed Group Relative Policy Optimization (GRPO) \citep{deepseekmath} further reduces memory requirements by estimating group scores, while replacing traditional reward models with rule-based reward mechanisms. Through this innovative paradigm, LLMs can significantly enhance their reasoning capabilities and domain-specific performance (e.g., in mathematics and programming) without requiring annotated data \citep{Deepseek-r1, o1}. Recent studies \citep{R1-searcher, rm-r1, search-r1, RecZero} have further explored the applications of RL in mathematics and general domains.

% \subsection{LLM post-training Integration}
% Reinforcement Learning (RL) \citep{RLsurvey, RLsurvey2} operates through the interaction mechanism between agents and the environment, rewarding correct actions with the core objective of maximizing cumulative returns. In recent years, researchers have introduced RL into the fine-tuning process of LLMs. Among these approaches, Reinforcement Learning with Human Feedback (RLHF) \citep{RLHF} typically employs the Proximal Policy Optimization (PPO) \citep{PPO} algorithm to achieve alignment with human preferences. However, the PPO algorithm faces technical challenges related to excessive memory consumption. To simplify the application of RL in LLMs, researchers have proposed innovative methods such as Direct Preference Optimization (DPO) \citep{DPO}. In the context of recommender systems, the s-DPO method \citep{sdpo} combines DPO with the softmax loss function, enhancing LLMs' capability in mining hard negative samples for recommender systems. Although these methods have improved efficiency, they still face challenges such as off-policy issues and performance ceilings that fall short of online methods. The recently proposed Group Relative Policy Optimization (GRPO) \citep{deepseekmath} further reduces memory requirements by estimating group scores, while replacing traditional reward models with rule-based reward mechanisms. Through this innovative paradigm, LLMs can significantly enhance their reasoning capabilities and domain-specific performance (e.g., in mathematics and programming) without requiring annotated data \citep{Deepseek-r1, o1}. Recent studies \citep{R1-searcher, rm-r1, search-r1} have further explored the applications of RL in mathematics and general domains. Despite these advancements, the application of RL in improving rating prediction performance through reasoning capabilities remains unexplored.

\section{Conclusion}
\label{conclusion}
This report introduces \textbf{MiniOneRec}, the first fully open-source generative recommendation framework, which to the best of our knowledge provides an end-to-end workflow spanning SID construction, SFT, and recommendation-oriented RL.
By systematically validating scaling laws on public benchmarks, we demonstrate that larger generative recommenders achieve lower training and evaluation losses than their smaller counterparts, confirming the parameter–efficiency advantage of the SID-based paradigm over traditional embedding-centric models.  
Building on this insight, we introduce post-training techniques:  
(i) full-process SID alignment, which embeds SID tokens into the model vocabulary and imposes auxiliary alignment tasks across both SFT and RL stages, and (ii) reinforced preference optimization, which combines constrained decoding, beam-based sampling, and hybrid reward design.  
Extensive experiments on Amazon Review show that MiniOneRec consistently surpasses strong sequential, generative, and LLM-based baselines, while maintaining a lean post-training footprint.

Looking forward, we will keep maintaining and extending the MiniOneRec codebase.  
A public roadmap will guide future developments, and we warmly welcome community contributions.  
Planned updates include new datasets, more advanced tokenisation schemes, larger backbone models, and enhanced training pipelines, ensuring that MiniOneRec remains a solid reference platform for research and practice in large-scale generative recommendation.



% \newpage

% \input{chapters/8_checklist}

\newpage
\bibliographystyle{unsrtnat}
\bibliography{iclr2026_conference}

\newpage
\appendix

\section{Alignment Prompts}
\label{appendix:align_prompt}

\subsection{Overview}
In this section we describe how we \textbf{align the world knowledge of a LLM with the SID space} so that the model can be directly used for generative recommendation.  

We append a three–layer codebook, each layer containing $256$ unique SIDs, to the original tokenizer vocabulary. The inserted codes are treated as indivisible tokens, enabling the LLM to read or write SID sequences without any sub-word splitting.
Then we introduce a set of auxiliary \emph{Alignment Tasks} to bridge the semantic gap between the textual vocabulary and the newly added SIDs. 
% These tasks repeatedly present the same item or user trace in two complementary views—plain text and SID sequences—so that the model is forced to learn a shared embedding space. 
A specialized \emph{Recommendation Tasks} further teach the model to predict the SID of the next item given an user history.

\subsection{Recommendation Tasks}

\begin{enumerate}
    \item \textbf{Generative Retrieval.} The main task ofor generative recommendation.  
    The LLM receives a chronologically ordered SID sequence that represents the user’s recent interactions, together with an explicit instruction such as “Recommend the next item.”, and is asked to predict the SID of the next item the user might engage with. A sample prompt is shown in the Figure \ref{fig:sid2sid}.

\begin{figure}[H]
    \centering
    \includegraphics[width=0.8\textwidth]{figures/sid2sid.png}
    \vspace{-10pt}
    \caption{Generative Retrieval Prompt.}
    \label{fig:sid2sid}
    \vspace{-10pt}
\end{figure}

    \item \textbf{Asymmetric Item Prediction.}
    (a) Given a textual user history, predict the SID of the next item; (b) given a SID-only history, generate the textual title of the next item. Sample prompts are shown in the Figure \ref{fig:titlehis2sid} and Figure \ref{fig:sidhis2title}.

\begin{figure}[H]
    \centering
    \includegraphics[width=0.8\textwidth]{figures/titlehis2sid.png}
    \vspace{-10pt}
    \caption{Asymmetric Item Prediction Prompt1.}
    \label{fig:titlehis2sid}
    \vspace{-10pt}
\end{figure}

\begin{figure}[H]
    \centering
    \includegraphics[width=0.8\textwidth]{figures/sidhis2title.png}
    \vspace{-10pt}
    \caption{Asymmetric Item Prediction Prompt2.}
    \label{fig:sidhis2title}
    \vspace{-10pt}
\end{figure}
\end{enumerate}

\subsection{Alignment Tasks}

\begin{enumerate}
    \item \textbf{SID–Text Semantic Alignment.} (a) Predict an item's textual title from its SID; (b) predict the SID from the item's textual title.  Sample prompts are shown in the Figure \ref{fig:sid2title} and Figure \ref{fig:title2sid}.

\begin{figure}[H]
    \centering
    \includegraphics[width=0.45\linewidth]{figures/sid2title.png}
    \vspace{-10pt}
    \caption{SID–Text Semantic Alignment Prompt1.}
    \label{fig:sid2title}
    \vspace{-10pt}
\end{figure}

\begin{figure}[H]
    \centering
    \includegraphics[width=0.64\textwidth]{figures/title2sid.png}
    \vspace{-10pt}
    \caption{SID–Text Semantic Alignment Prompt2.}
    \label{fig:title2sid}
    \vspace{-10pt}
\end{figure}

    \item \textbf{Item Description Reconstruction.}  
          To ground SIDs in richer semantics, we ask the model to generate the item description from a single SID and, conversely, infer the SID from the description. We perform this task only during the SFT stage because the description space is large and diverse. Sample prompts is shown in the Figure \ref{fig:description2sid}.

\begin{figure}[H]
    \centering
    \includegraphics[width=0.8\textwidth]{figures/description2sid.png}
    \vspace{-10pt}
    \caption{Item Description Reconstruction Prompt.}
    \label{fig:description2sid}
    \vspace{-10pt}
\end{figure}

    \item \textbf{User Preference Summarization.}  
          Given a sequence of SIDs, the model produces a short natural–language profile that summarizes the user’s interests. Because the raw dataset lacks explicit preference labels, we employ \textsc{DeepSeek} \citep{deepseek} to extract summaries from the item's meta data and users’ textual reviews and use them as pseudo labels. This task, too, is restricted to the SFT stage due to the open-ended output space.
\begin{figure}[H]
    \centering
    \includegraphics[width=0.8\textwidth]{figures/sidhis2preference.png}
    \vspace{-10pt}
    \caption{User Preference Summarization Prompt.}
    \label{fig:sid2title}
    \vspace{-10pt}
\end{figure}

\end{enumerate}


\section{Datasets}

We evaluate our model on two Amazon Review subsets:  Amazon Review subsets (\textit{Industrial\_and\_Scientific} and \textit{Office\_Products}).  
To keep computational costs affordable, we follow the trimming strategy used in \citep{D3}.  
The main steps are:  
(1) remove users and items with fewer than five interactions;  
(2) for \textit{Toys\_and\_Games}, keep events from October 2016 to November 2018;  
(3) for \textit{Industrial\_and\_Scientific}, which is smaller, keep all events between October 1996 and November 2018;   
(4) truncate every user history to at most ten items;  
(5) finally, split each dataset chronologically into training, validation and test sets with an 8:1:1 ratio.  
Key statistics of the resulting training splits are reported in Table~\ref{tab:stats}.

\begin{table}[H]
\centering
\caption{Statistics of datasets.} 
\label{tab:stats}
\begin{tabular}{ccc}
\toprule
Datasets & Inductrial & Office \\ \midrule
Items    & 3,685      & 3,459  \\
Train    & 3,6259     & 3,8924 \\
Valid    & 4,532      & 4,866  \\
Test     & 4,533      & 4,866  \\ \bottomrule
\end{tabular}
\end{table}


\end{document}
